\chapter{Special Functions}\label{special-functions}
\chapterstatus{draft}

The functions of classical analysis that are not elementary but appear everywhere --- in number theory,
probability, mathematical physics and geometry --- share a common theory: they are solutions of linear
differential equations with few singular points, they have integral representations, and their
asymptotics can be computed by a handful of methods. This chapter treats them together.
\begin{itemize}
\item Section~\ref{special-functions:section-gamma}: the Gamma and Beta functions, Euler's limit formula and the Weierstrass product, the cotangent and sine expansions, the reflection and
duplication formulas, and the Bohr--Mollerup theorem.
\item Section~\ref{special-functions:section-asymptotics}: asymptotic expansions, Watson's lemma, Laplace's method and Stirling's formula, stationary phase, Bernoulli polynomials and Euler--Maclaurin summation.
\item Section~\ref{special-functions:section-orthogonal-polynomials}: orthogonal polynomials, the three-term recurrence, zeros, Gauss quadrature, the classical families and Chebyshev's extremal property.
\item Section~\ref{special-functions:section-fuchsian}: linear differential equations in the complex domain, monodromy, regular singular points and the Frobenius method.
\item Section~\ref{special-functions:section-hypergeometric}: the hypergeometric function, its equation, Euler's integral, Gauss's summation theorem and the Chu--Vandermonde identity.
\item Section~\ref{special-functions:section-bessel}: Bessel functions and spherical harmonics.
\item Section~\ref{special-functions:section-elliptic}: elliptic integrals, the arithmetic--geometric mean, and Weierstrass's $\wp$-function.
\item Section~\ref{special-functions:section-theta}: Jacobi's theta function, Poisson summation and the transformation formula.
\end{itemize}

\noindent\textbf{Prerequisites.} Chapter~\ref{real-analysis} (Real Analysis), Chapter~\ref{linear-algebra} (Linear Algebra), Chapter~\ref{generating-functions} (Generating Functions),
Chapter~\ref{measure-theory} (Measure Theory) and Chapter~\ref{complex-analysis} (Complex Analysis).

\section{The Gamma function}\label{special-functions:section-gamma}

Euler's Gamma function interpolates the factorials. We develop its basic properties: meromorphic continuation, Euler's limit formula and the Weierstrass product, absence of zeros, the
reflection and duplication formulas, and the characterisation by logarithmic convexity. Along the way we prove the partial fraction expansion of the cotangent and the product formula for the sine,
which also give the values $\zeta(2k)$ in Chapter~\ref{analytic-number-theory}. Integrals depending holomorphically on a parameter are handled by the following lemma.

\begin{lemma}\label{special-functions:lemma-holomorphic-integrals}
Let $I \subseteq \RR$ be an interval, $U \subseteq \CC$ open, and $F \colon I \times U \to \CC$ continuous and holomorphic in $s$ for each fixed $x \in I$. Suppose that for every compact $K \subseteq U$
there is an integrable function $G_K$ on $I$ with $|F(x, s)| \leq G_K(x)$ for $x \in I$, $s \in K$. Then $\Phi(s) = \int_IF(x, s)\,dx$ is holomorphic on $U$, and
$\Phi^{\prime}(s) = \int_I\partial_sF(x, s)\,dx$.
\end{lemma}

\begin{proof}
\emph{Continuity.} If $s_k \to s$ in $U$, the points lie in a compact disc $K \subseteq U$ for large $k$, $F(x, s_k) \to F(x, s)$ for each $x$, and $|F(x, s_k)| \leq G_K(x)$; so
$\Phi(s_k) \to \Phi(s)$ by dominated convergence (Theorem~\ref{measure-theory:theorem-dominated-convergence}).
\emph{Holomorphy.} Let $T \subseteq U$ be a closed triangle. The function $(x, w) \mapsto F(x, w)$ on $I \times \partial T$ is continuous and bounded by $G_T(x)$, so by Fubini's theorem
(Theorem~\ref{measure-theory:theorem-fubini}) $\int_{\partial T}\Phi(w)\,dw = \int_I\int_{\partial T}F(x, w)\,dw\,dx = 0$, the inner integral vanishing by Goursat's theorem
(Theorem~\ref{complex-analysis:theorem-goursat}). By Morera's theorem (Corollary~\ref{complex-analysis:corollary-liouville}(3)), $\Phi$ is holomorphic.
\emph{Derivative.} Let $\overline{D}(s, r) \subseteq U$. By the Cauchy integral formula for derivatives (Theorem~\ref{complex-analysis:theorem-cauchy-formula}), applied to $\Phi$ and to each
$F(x, \cdot)$, and Fubini's theorem again on $I \times \partial D(s, r)$,
$\Phi^{\prime}(s) = \frac{1}{2\pi i}\int_{\partial D}\frac{\Phi(w)}{(w - s)^2}dw = \int_I\frac{1}{2\pi i}\int_{\partial D}\frac{F(x, w)}{(w - s)^2}dw\,dx = \int_I\partial_sF(x, s)\,dx$.
\end{proof}

\begin{definition}\label{special-functions:definition-gamma}
The \emph{Gamma function} is
\[
\Gamma(s) = \int_0^\infty t^{s-1}e^{-t}\,dt, \qquad \sigma > 0,
\]
where $t^{s-1} = e^{(s-1)\log t}$.
\end{definition}

\begin{proposition}\label{special-functions:proposition-gamma-basic}
\begin{enumerate}
\item The integral converges absolutely for $\sigma > 0$ and defines a holomorphic function there, with $\Gamma(\sigma) > 0$ for real $\sigma > 0$ and $\Gamma(\bar s) = \overline{\Gamma(s)}$.
\item $\Gamma(s + 1) = s\Gamma(s)$ for $\sigma > 0$, $\Gamma(1) = 1$, and $\Gamma(n + 1) = n!$ for integers $n \geq 0$.
\item $\Gamma$ extends to a meromorphic function on $\CC$, holomorphic except for simple poles at $s = 0, -1, -2, \ldots$, with $\operatorname{Res}_{s = -n}\Gamma = \frac{(-1)^n}{n!}$, and
$\Gamma(s + 1) = s\Gamma(s)$ holds on all of $\CC$.
\end{enumerate}
\end{proposition}

\begin{proof}
(1) If $a \leq \sigma \leq b$ with $0 < a$, then $|t^{s-1}e^{-t}| = t^{\sigma - 1}e^{-t} \leq (t^{a-1} + t^{b-1})e^{-t}$, which is integrable on $(0, \infty)$ (near $0$ because $a > 0$, near $\infty$
because of $e^{-t}$). So Lemma~\ref{special-functions:lemma-holomorphic-integrals} applies with $I = (0, \infty)$. For real $s$ the integrand is positive, and conjugating the integrand gives
$\Gamma(\bar s) = \overline{\Gamma(s)}$.
(2) Integrate by parts on $[\varepsilon, R]$: $\int_\varepsilon^Rt^se^{-t}dt = [-t^se^{-t}]_\varepsilon^R + s\int_\varepsilon^Rt^{s-1}e^{-t}dt$, and let $\varepsilon \to 0$, $R \to \infty$. Also
$\Gamma(1) = \int_0^\infty e^{-t}dt = 1$, and induction gives $\Gamma(n + 1) = n!$.
(3) For an integer $n \geq 1$ put $\Gamma_n(s) = \frac{\Gamma(s + n)}{s(s + 1)\cdots(s + n - 1)}$ on $\sigma > -n$, $s \notin \{0, -1, \ldots, -n + 1\}$. It is meromorphic there and, by (2),
equals $\Gamma$ on $\sigma > 0$ and $\Gamma_{n+1}$ on $\sigma > -n$. So these functions glue to a meromorphic function on $\CC$. Near $s = -n$, $\Gamma_{n+1}$ is $\Gamma(s + n + 1)$, holomorphic and
equal to $1$ at $s = -n$, divided by $s(s + 1)\cdots(s + n)$, which has a simple zero at $-n$; so the pole is simple with residue
$\frac{1}{(-n)(-n + 1)\cdots(-1)} = \frac{(-1)^n}{n!}$. The functional equation persists by the identity theorem (Theorem~\ref{complex-analysis:theorem-identity}).
\end{proof}

The Gamma function is related to many integrals through the following formula, which uses a change of variables in a double integral.

\begin{proposition}[Euler's Beta integral]\label{special-functions:proposition-beta}
For $\operatorname{Re}a, \operatorname{Re}b > 0$, let $B(a, b) = \int_0^1u^{a-1}(1 - u)^{b-1}\,du$. Then
\[
\Gamma(a)\Gamma(b) = \Gamma(a + b)B(a, b), \qquad \text{so} \quad B(a, b) = \frac{\Gamma(a)\Gamma(b)}{\Gamma(a + b)}
\]
whenever $\Gamma(a + b) \neq 0$: for real $a, b > 0$ by Proposition~\ref{special-functions:proposition-gamma-basic}(1), and always by Theorem~\ref{special-functions:theorem-weierstrass-gamma} below.
\end{proposition}

\begin{proof}
By Fubini's theorem (Theorem~\ref{measure-theory:theorem-fubini}), the integrand being absolutely integrable,
$\Gamma(a)\Gamma(b) = \int_0^\infty\int_0^\infty x^{a-1}y^{b-1}e^{-x-y}\,dx\,dy$. The map $(r, u) \mapsto (x, y) = (ru, r(1 - u))$ is a diffeomorphism from $(0, \infty) \times (0, 1)$ onto
$(0, \infty)^2$, with inverse $r = x + y$, $u = x/(x + y)$, and Jacobian determinant $\bigl|\det\begin{pmatrix}u & r \\ 1 - u & -r\end{pmatrix}\bigr| = r$. By the change of variables formula
(Theorem~\ref{measure-theory:theorem-change-of-variables}),
\[
\Gamma(a)\Gamma(b) = \int_0^\infty\int_0^1r^{a+b-1}u^{a-1}(1 - u)^{b-1}e^{-r}\,du\,dr = \Gamma(a + b)B(a, b). \qedhere
\]
\end{proof}

\begin{example}\label{special-functions:example-gamma-half}
With $a = b = \frac12$ and the substitution $u = \sin^2\phi$, $du = 2\sin\phi\cos\phi\,d\phi$,
\[
\Gamma(\tfrac12)^2 = \Gamma(1)B(\tfrac12, \tfrac12) = \int_0^1\frac{du}{\sqrt{u(1 - u)}} = \int_0^{\pi/2}2\,d\phi = \pi, \qquad \text{so} \quad \Gamma(\tfrac12) = \sqrt{\pi}
\]
(as $\Gamma(\frac12) > 0$). Substituting $t = x^2$ in the definition gives $\Gamma(\frac12) = \int_0^\infty e^{-x^2}\cdot 2\,dx$, that is, the Gaussian integral $\int_{-\infty}^\infty e^{-x^2}dx = \sqrt{\pi}$. Then
$\Gamma(\frac32) = \frac12\sqrt{\pi}$, $\Gamma(\frac52) = \frac34\sqrt{\pi}$, and the volume of the unit ball in $\RR^n$ is $\omega_n = \pi^{n/2}/\Gamma(\frac n2 + 1)$
(Exercise~\ref{measure-theory:exercise-ball-volume}): $\omega_3 = \pi^{3/2}/\Gamma(\frac52) = \frac43\pi$.
\end{example}

\begin{theorem}[Euler's limit formula]\label{special-functions:theorem-euler-limit}
For $\sigma > 0$,
\[
\Gamma(s) = \lim_{n \to \infty}\frac{n!\,n^s}{s(s + 1)\cdots(s + n)}.
\]
\end{theorem}

\begin{proof}
\emph{Step 1: a truncated integral.} Let $\Gamma_n(s) = \int_0^nt^{s-1}(1 - t/n)^n\,dt$. Substituting $t = nu$ gives $\Gamma_n(s) = n^sB(s, n + 1)$. For $m \geq 1$, integration by parts gives
$\int_0^1u^{a-1}(1 - u)^m\,du = \frac ma\int_0^1u^a(1 - u)^{m-1}\,du$ when $\operatorname{Re}a > 0$ (the boundary terms vanish), that is $B(a, m + 1) = \frac maB(a + 1, m)$. Applying this $n$ times,
\[
\Gamma_n(s) = n^s\frac{n!}{s(s + 1)\cdots(s + n - 1)}\int_0^1u^{s+n-1}\,du = \frac{n!\,n^s}{s(s + 1)\cdots(s + n)}.
\]
\emph{Step 2: the limit.} Since $\log(1 - v) \leq -v$ for $0 \leq v < 1$, $0 \leq (1 - t/n)^n\mathbf{1}_{t < n} \leq e^{-t}$, and $(1 - t/n)^n \to e^{-t}$ for each $t$
(Proposition~\ref{real-analysis:proposition-notable-limits}). By dominated convergence, with dominating function $t^{\sigma - 1}e^{-t}$, $\Gamma_n(s) \to \Gamma(s)$.
\end{proof}

To pass from the limit formula to a product we need infinite products of holomorphic functions.

\begin{lemma}[Infinite products]\label{special-functions:lemma-infinite-products}
Let $U \subseteq \CC$ be open and $u_k$ holomorphic functions on $U$ such that $\sum_k\sup_K|u_k| < \infty$ for every compact $K \subseteq U$. Then the partial products $P_N = \prod_{k \leq N}(1 + u_k)$
converge uniformly on compact subsets of $U$ to a holomorphic function $P$. A point $s$ is a zero of $P$ exactly when $1 + u_k(s) = 0$ for some $k$, and its order is the sum of the orders of these
zeros. Off the zeros, $P^{\prime}/P = \sum_ku_k^{\prime}/(1 + u_k)$, uniformly on compact subsets.
\end{lemma}

\begin{proof}
\emph{Step 1: a logarithm.} For $|u| < 1$ let $\log(1 + u) = \sum_{j \geq 1}(-1)^{j+1}u^j/j$. Its derivative is $\sum_{j \geq 0}(-u)^j = (1 + u)^{-1}$, so
$\frac{d}{du}\bigl((1 + u)e^{-\log(1 + u)}\bigr) = 0$, and $\exp(\log(1 + u)) = 1 + u$ because both sides agree at $u = 0$. For $|u| \leq \frac12$, and $|\log(1 + u)| \leq \sum_j|u|^j = \frac{|u|}{1 - |u|} \leq 2|u|$.
\emph{Step 2: convergence.} Fix a compact $K$. There is $k_0$ with $|u_k| \leq \frac12$ on $K$ for $k > k_0$. The series $\sum_{k > k_0}\log(1 + u_k)$ converges uniformly on $K$ by Step 1 and the hypothesis
(Weierstrass M-test), to a continuous function $S$, so $\prod_{k_0 < k \leq N}(1 + u_k) = \exp(\sum_{k_0 < k \leq N}\log(1 + u_k)) \to e^S$ uniformly on $K$, and $e^S$ has no zeros. Multiplying by the
finitely many factors $k \leq k_0$ gives the uniform convergence of $P_N$ on $K$, and the zeros of $P$ in $K$ are those of these finitely many factors, with orders adding up. $P$ is holomorphic by
Weierstrass' theorem (Corollary~\ref{complex-analysis:corollary-liouville}(4)).
\emph{Step 3: logarithmic derivative.} $P_N^{\prime}/P_N = \sum_{k \leq N}u_k^{\prime}/(1 + u_k)$ by the product rule, and $P_N \to P$, $P_N^{\prime} \to P^{\prime}$ uniformly on compact subsets
(Weierstrass' theorem); on a compact set where $P \neq 0$, $|P|$ is bounded below, so $P_N^{\prime}/P_N \to P^{\prime}/P$ uniformly there.
\end{proof}

\begin{lemma}\label{special-functions:lemma-euler-constant}
The limit $\gamma = \lim_{n \to \infty}(H_n - \log n)$, where $H_n = 1 + \frac12 + \cdots + \frac1n$, exists, and $0 < \gamma < 1$. It is \emph{Euler's constant}, $\gamma = 0.5772156649\ldots$
\end{lemma}

\begin{proof}
Let $a_n = H_n - \log n$ and $d_k = a_k - a_{k+1} = \int_k^{k+1}\frac{dx}{x} - \frac{1}{k + 1}$. Since $\frac{1}{k+1} < \frac1x < \frac1k$ on $(k, k + 1)$, we have
$0 < d_k < \frac1k - \frac{1}{k + 1}$. So $(a_n)$ is decreasing, and $a_n = a_1 - \sum_{k<n}d_k > 1 - \sum_{k<n}\bigl(\frac1k - \frac{1}{k+1}\bigr) = \frac1n > 0$. Hence $a_n$ converges, and
$\gamma = 1 - \sum_{k \geq 1}d_k$ with $0 < \sum_kd_k < \sum_k\bigl(\frac1k - \frac{1}{k+1}\bigr) = 1$, so $0 < \gamma < 1$.
\end{proof}

\begin{theorem}[Weierstrass product]\label{special-functions:theorem-weierstrass-gamma}
The function $1/\Gamma$ extends to an entire function, and for all $s \in \CC$
\[
\frac{1}{\Gamma(s)} = se^{\gamma s}\prod_{k \geq 1}\Bigl(1 + \frac sk\Bigr)e^{-s/k},
\]
the product converging uniformly on compact sets. Its zeros are the points $0, -1, -2, \ldots$, all simple. In particular $\Gamma$ has no zeros, and Euler's limit formula
(Theorem~\ref{special-functions:theorem-euler-limit}) holds for all $s \notin \{0, -1, -2, \ldots\}$.
\end{theorem}

\begin{proof}
\emph{Step 1: the product converges.} For $|w| \leq 1$, $(1 + w)e^{-w} = 1 + \sum_{j \geq 2}c_jw^j$ with $c_j = \frac{(-1)^j}{j!} + \frac{(-1)^{j-1}}{(j - 1)!} = (-1)^j\frac{1 - j}{j!}$, so
$|(1 + w)e^{-w} - 1| \leq |w|^2\sum_{j \geq 2}\frac{j - 1}{j!} = |w|^2$ (the last sum is $\sum_{j \geq 1}\frac{1}{(j - 1)!} - \sum_{j \geq 1}\frac1{j!} = e - (e - 1) = 1$, the term $j = 1$ being $0$). With
$w = s/k$, the factors $1 + u_k(s) = (1 + s/k)e^{-s/k}$ satisfy $|u_k(s)| \leq R^2/k^2$ for $|s| \leq R \leq k$. By Lemma~\ref{special-functions:lemma-infinite-products}, $E(s) = se^{\gamma s}\prod_k(1 + s/k)e^{-s/k}$
is an entire function, with simple zeros exactly at $0, -1, -2, \ldots$
\emph{Step 2: $E = 1/\Gamma$.} For $\sigma > 0$, with $H_n = \sum_{k \leq n}\frac1k$,
\[
\frac{s(s + 1)\cdots(s + n)}{n!\,n^s} = s\prod_{k=1}^n\Bigl(1 + \frac sk\Bigr)e^{-s\log n} = se^{s(H_n - \log n)}\prod_{k=1}^n\Bigl(1 + \frac sk\Bigr)e^{-s/k},
\]
which tends to $E(s)$, because $H_n - \log n \to \gamma$ (Lemma~\ref{special-functions:lemma-euler-constant}). By Theorem~\ref{special-functions:theorem-euler-limit} the left side
tends to $1/\Gamma(s)$; so $\Gamma(s)E(s) = 1$ for $\sigma > 0$. Both sides are meromorphic on the connected set $\CC$, so $\Gamma E = 1$ everywhere by the identity theorem. Hence $1/\Gamma = E$ is
entire, and $\Gamma$ never vanishes. Finally, the computation of Step 2 shows that $\frac{s(s + 1)\cdots(s + n)}{n!\,n^s} \to E(s)$ for every $s \in \CC$; for $s \notin \{0, -1, \ldots\}$ the limit
$E(s) = 1/\Gamma(s)$ is nonzero, and taking reciprocals gives Euler's limit formula at $s$.
\end{proof}

\subsection{The cotangent, the sine, and the reflection formula}\label{special-functions:subsection-cotangent}

\begin{theorem}[Euler]\label{special-functions:theorem-cotangent}
For $z \in \CC \setminus \ZZ$,
\[
\pi\cot\pi z = \frac1z + \sum_{n \geq 1}\frac{2z}{z^2 - n^2} = \lim_{N \to \infty}\sum_{n = -N}^N\frac{1}{z - n},
\]
the series converging absolutely and uniformly on compact subsets of $\CC \setminus \ZZ$.
\end{theorem}

\begin{proof}
We follow Herglotz's argument \cite{aigner-ziegler}.
\emph{Step 1: the series.} For $|z| \leq R$ and $n \geq 2R$, $|\frac{2z}{z^2 - n^2}| \leq \frac{2R}{n^2 - R^2} \leq \frac{8R}{3n^2}$, so the series converges absolutely and locally uniformly on
$\CC \setminus \ZZ$; call its sum $S(z)$. Since $\frac{1}{z - n} + \frac{1}{z + n} = \frac{2z}{z^2 - n^2}$, $S(z) = \lim_NS_N(z)$ with $S_N(z) = \sum_{|n| \leq N}\frac{1}{z - n}$. Near an integer $m$,
$S(z) - \frac{1}{z - m}$ is a locally uniform limit of functions holomorphic near $m$; so $S$ has a simple pole with residue $1$ at each integer.
\emph{Step 2: the difference is entire.} $\pi\cot\pi z = \pi\cos\pi z/\sin\pi z$ also has simple poles with residue $\frac{\pi\cos\pi m}{\pi\cos\pi m} = 1$ at each integer $m$, and no other poles. So
$h(z) = \pi\cot\pi z - S(z)$ has removable singularities at the integers (Proposition~\ref{complex-analysis:proposition-singularities}) and is entire. Both functions are odd, so $h$ is odd and
$h(0) = 0$; both are real on $\RR \setminus \ZZ$, so $h$ is real on $\RR$.
\emph{Step 3: periodicity.} $S_N(z + 1) - S_N(z) = \frac{1}{z + 1 + N} - \frac{1}{z - N} \to 0$, so $S$ has period $1$, as $\cot\pi z$ does; hence $h(z + 1) = h(z)$.
\emph{Step 4: a duplication identity.} For $f = \pi\cot\pi z$: with $a = \pi z/2$, $\cot a + \cot(a + \frac\pi2) = \cot a - \tan a = \frac{\cos^2a - \sin^2a}{\sin a\cos a} = 2\cot 2a$, so
$f(\frac z2) + f(\frac{z + 1}2) = 2f(z)$. For $S$: $S_N(\frac z2) + S_N(\frac{z + 1}2) = \sum_{|n| \leq N}\bigl(\frac{2}{z - 2n} + \frac{2}{z - (2n - 1)}\bigr) = 2S_{2N}(z) + \frac{2}{z + 2N + 1}$, since the
numbers $2n$ and $2n - 1$ with $|n| \leq N$ are exactly the integers from $-2N - 1$ to $2N$. Letting $N \to \infty$, $S(\frac z2) + S(\frac{z + 1}2) = 2S(z)$. Hence $h(\frac z2) + h(\frac{z + 1}2) = 2h(z)$
for $z \notin \ZZ$, and by continuity for all $z$.
\emph{Step 5: Herglotz's trick.} The continuous periodic real function $h|_\RR$ attains its maximum $M$ at some $x_0 \in [0, 1]$. Then $h(\frac{x_0}2) + h(\frac{x_0 + 1}2) = 2M$ with both terms $\leq M$, so
$h(\frac{x_0}2) = M$. Repeating, $h(x_0/2^k) = M$ for all $k$, and by continuity $M = h(0) = 0$. Applying this to $-h$, the minimum is $0$ too, so $h = 0$ on $\RR$, and $h = 0$ on $\CC$ by the identity
theorem.
\end{proof}

\begin{corollary}[Euler's product for the sine]\label{special-functions:corollary-sine-product}
For all $z \in \CC$,
\[
\sin\pi z = \pi z\prod_{n \geq 1}\Bigl(1 - \frac{z^2}{n^2}\Bigr).
\]
\end{corollary}

\begin{proof}
By Lemma~\ref{special-functions:lemma-infinite-products} with $u_n = -z^2/n^2$, the right side is an entire function $P$ with simple zeros exactly at the integers, and off the integers
$\frac{P^{\prime}}{P} = \frac1z + \sum_n\frac{-2z/n^2}{1 - z^2/n^2} = \frac1z + \sum_n\frac{2z}{z^2 - n^2} = \pi\cot\pi z = \frac{(\sin\pi z)^{\prime}}{\sin\pi z}$ by Theorem~\ref{special-functions:theorem-cotangent}.
So $(\sin\pi z/P)^{\prime} = \frac{\sin\pi z}{P}\bigl(\frac{(\sin\pi z)^{\prime}}{\sin\pi z} - \frac{P^{\prime}}{P}\bigr) = 0$ on the connected open set $\CC \setminus \ZZ$, and $\sin\pi z = cP(z)$. Letting $z \to 0$,
$c = \lim\frac{\sin\pi z}{\pi z} = 1$.
\end{proof}

\begin{example}\label{special-functions:example-wallis}
At $z = \frac12$: $1 = \frac\pi2\prod_n\bigl(1 - \frac{1}{4n^2}\bigr)$, that is, \emph{Wallis' product} $\frac\pi2 = \prod_{n \geq 1}\frac{(2n)^2}{(2n - 1)(2n + 1)} = \frac21\cdot\frac23\cdot\frac43\cdot\frac45\cdot\frac65\cdot\frac67\cdots$
Comparing the coefficients of $z^3$ in $\sin\pi z = \pi z - \frac{\pi^3z^3}{6} + \cdots$ and in the product, $-\frac{\pi^3}{6} = -\pi\sum_n\frac1{n^2}$, which is Euler's original (heuristic) derivation of
$\zeta(2) = \frac{\pi^2}{6}$; here it is a proof.
\end{example}

\begin{theorem}[Reflection formula]\label{special-functions:theorem-reflection}
For $s \in \CC \setminus \ZZ$,
\[
\Gamma(s)\Gamma(1 - s) = \frac{\pi}{\sin\pi s}.
\]
\end{theorem}

\begin{proof}
By Theorem~\ref{special-functions:theorem-weierstrass-gamma}, the factors $e^{\pm\gamma s}$ and $e^{\mp s/k}$ cancel in
$\frac{1}{\Gamma(s)\Gamma(-s)} = se^{\gamma s}\prod_k\bigl(1 + \frac sk\bigr)e^{-s/k}\cdot(-s)e^{-\gamma s}\prod_k\bigl(1 - \frac sk\bigr)e^{s/k} = -s^2\prod_k\Bigl(1 - \frac{s^2}{k^2}\Bigr) = -\frac{s\sin\pi s}{\pi}$
by Corollary~\ref{special-functions:corollary-sine-product}. Since $\Gamma(1 - s) = -s\Gamma(-s)$, $\Gamma(s)\Gamma(1 - s) = -s\Gamma(s)\Gamma(-s) = \frac{\pi}{\sin\pi s}$.
\end{proof}

\begin{example}\label{special-functions:example-reflection}
At $s = \frac12$ the formula gives $\Gamma(\frac12)^2 = \pi$ again. On the line $s = \frac12 + it$, since $\Gamma(\frac12 - it) = \overline{\Gamma(\frac12 + it)}$ and
$\sin(\frac\pi2 + i\pi t) = \cos(i\pi t) = \cosh\pi t$,
\[
\bigl|\Gamma(\tfrac12 + it)\bigr|^2 = \frac{\pi}{\cosh\pi t}:
\]
the Gamma function decays exponentially on vertical lines. This decay, of size $e^{-\pi|t|/2}$ on every vertical line by Stirling's formula (Remark~\ref{special-functions:remark-stirling-gamma}), is what
makes the contour integrals of Chapter~\ref{analytic-number-theory} converge.
\end{example}

\begin{proposition}[Legendre's duplication formula]\label{special-functions:proposition-duplication}
For $2s \notin \{0, -1, -2, \ldots\}$, $\Gamma(s)\Gamma(s + \frac12) = 2^{1 - 2s}\sqrt{\pi}\,\Gamma(2s)$.
\end{proposition}

\begin{proof}
Let $s > 0$ be real. In $B(s, s) = \int_0^1(u(1 - u))^{s-1}du = 2\int_0^{1/2}(u(1 - u))^{s-1}du$ substitute $v = 4u(1 - u)$, which increases from $0$ to $1$ on $[0, \frac12]$, with
$dv = 4(1 - 2u)\,du$ and $1 - 2u = \sqrt{1 - v}$:
\[
B(s, s) = 2\int_0^1\Bigl(\frac v4\Bigr)^{s-1}\frac{dv}{4\sqrt{1 - v}} = 2^{1 - 2s}B(s, \tfrac12).
\]
By Proposition~\ref{special-functions:proposition-beta}, $\frac{\Gamma(s)^2}{\Gamma(2s)} = 2^{1-2s}\frac{\Gamma(s)\Gamma(\frac12)}{\Gamma(s + \frac12)}$, which is the claim for real $s > 0$, as
$\Gamma(\frac12) = \sqrt{\pi}$. Both sides are meromorphic on $\CC$, so the identity theorem extends it.
\end{proof}

\begin{remark}\label{special-functions:remark-stirling-gamma}
Stirling's formula for $n!$ (Theorem~\ref{real-analysis:theorem-stirling}) extends to the complex Gamma function: for every $\delta > 0$,
$\log\Gamma(s) = (s - \frac12)\log s - s + \frac12\log 2\pi + O_\delta(|s|^{-1})$ uniformly for $|\arg s| \leq \pi - \delta$, $|s| \geq 1$, with the principal branches of the logarithms
\cite{iwaniec-kowalski}, \cite{titchmarsh-zeta}. In particular $|\Gamma(\sigma + it)| \sim \sqrt{2\pi}|t|^{\sigma - 1/2}e^{-\pi|t|/2}$ as $|t| \to \infty$ with $\sigma$ fixed, in agreement with
Example~\ref{special-functions:example-reflection}. For real $s \to \infty$ the leading term is proved in Corollary~\ref{special-functions:corollary-stirling-gamma} by Laplace's method.
\end{remark}

\subsection{Logarithmic convexity}\label{special-functions:subsection-log-convexity}

The functional equation $f(x + 1) = xf(x)$ with $f(1) = 1$ does not determine $\Gamma$: one may multiply by any function of period $1$ equal to $1$ at the integers, such as $1 + \sin^2\pi x$. Bohr and
Mollerup found a natural extra condition.

\begin{definition}\label{special-functions:definition-log-convex}
A function $f \colon (0, \infty) \to (0, \infty)$ is \emph{logarithmically convex} if $\log f$ is convex.
\end{definition}

\begin{theorem}[Bohr--Mollerup]\label{special-functions:theorem-bohr-mollerup}
$\Gamma$ is the only logarithmically convex function $f \colon (0, \infty) \to (0, \infty)$ with $f(x + 1) = xf(x)$ and $f(1) = 1$.
\end{theorem}

\begin{proof}
\emph{$\Gamma$ is logarithmically convex.} For $x, y > 0$ and $0 < \lambda < 1$, H\"older's inequality (Proposition~\ref{measure-theory:proposition-holder}) with exponents $1/\lambda$ and $1/(1 - \lambda)$ gives
$\Gamma(\lambda x + (1 - \lambda)y) = \int_0^\infty(t^{x-1}e^{-t})^\lambda(t^{y-1}e^{-t})^{1-\lambda}dt \leq \Gamma(x)^\lambda\Gamma(y)^{1-\lambda}$.
\emph{Uniqueness.} Let $f$ be as in the statement and $\varphi = \log f$; then $\varphi(n + 1) = \log n!$. Let $0 < x \leq 1$ and $n \geq 2$. For a convex function, slopes of chords increase, so comparing the chords on
$[n - 1, n]$, $[n, n + x]$ and $[n, n + 1]$,
\[
\log(n - 1) = \varphi(n) - \varphi(n - 1) \leq \frac{\varphi(n + x) - \varphi(n)}{x} \leq \varphi(n + 1) - \varphi(n) = \log n.
\]
By the functional equation $\varphi(n + x) = \varphi(x) + \log\bigl(x(x + 1)\cdots(x + n - 1)\bigr)$, so
\[
0 \leq \varphi(x) - \log\frac{(n - 1)!\,(n - 1)^x}{x(x + 1)\cdots(x + n - 1)} \leq x\log\frac{n}{n - 1}.
\]
Letting $n \to \infty$, $\varphi(x)$ is the limit of a quantity that does not depend on $f$. So $f$ is determined on $(0, 1]$, and by the functional equation everywhere; as $\Gamma$ has the properties, $f = \Gamma$.
\end{proof}

\begin{remark}\label{special-functions:remark-bohr-mollerup}
The proof recovers Euler's limit formula (Theorem~\ref{special-functions:theorem-euler-limit}) for real $s$, and the method gives short proofs of the duplication formula: $2^{x-1}\Gamma(\frac x2)\Gamma(\frac{x+1}2)/\sqrt{\pi}$ is
logarithmically convex, satisfies the functional equation and equals $1$ at $x = 1$, so it is $\Gamma(x)$ (Exercise~\ref{special-functions:exercise-duplication-bohr-mollerup}). Without logarithmic convexity uniqueness fails: for $\varepsilon > 0$,
$f_\varepsilon(x) = \Gamma(x)(1 + \varepsilon\sin^2\pi x)$ satisfies the functional equation and $f_\varepsilon(1) = 1$. It is not logarithmically convex: at the points $x = n + \frac12$ the second derivative of
$\log(1 + \varepsilon\sin^2\pi x)$ equals $-\frac{2\pi^2\varepsilon}{1 + \varepsilon}$, while $(\log\Gamma)^{\prime\prime}(x) = \sum_{k \geq 0}(x + k)^{-2}$ (Exercise~\ref{special-functions:exercise-digamma}) tends to $0$, so
$(\log f_\varepsilon)^{\prime\prime} < 0$ at $n + \frac12$ for large $n$.
\end{remark}

\section{Asymptotic methods}\label{special-functions:section-asymptotics}

Most special functions are given by integrals or series whose behaviour for large values of a parameter is what one needs in applications. This section develops the basic methods.

\begin{definition}\label{special-functions:definition-asymptotic}
Let $f$ be defined for large $x$ (real, or complex in a sector). A formal series $\sum_{k \geq 0}a_kx^{-k}$ is an \emph{asymptotic expansion} of $f$ as $x \to \infty$, written $f(x) \sim \sum_ka_kx^{-k}$, if for every
$N \geq 0$,
\[
f(x) - \sum_{k=0}^{N-1}a_kx^{-k} = O(x^{-N}) \qquad (x \to \infty).
\]
More generally one allows $f \sim \phi(x)\sum_ka_kx^{-k-\alpha}$ for a fixed function $\phi$ and exponent $\alpha$. For functions, $f \sim g$ means $f/g \to 1$.
\end{definition}

\begin{notation}\label{special-functions:notation-asymptotic}
$f = O(g)$ means $|f| \leq C|g|$ for large $x$; $f = o(g)$ means $f/g \to 0$; $f \sim g$ means $f/g \to 1$. The symbol $\sim$ between a function and a series refers to
Definition~\ref{special-functions:definition-asymptotic}.
\end{notation}

The coefficients of an asymptotic expansion are unique, $a_N = \lim_xx^N(f(x) - \sum_{k<N}a_kx^{-k})$, but the series need not converge, and when it converges its sum need not be $f$.

\begin{example}\label{special-functions:example-stieltjes}
Let $f(x) = \int_0^\infty\frac{e^{-xt}}{1 + t}\,dt$ for $x > 0$. Expanding $\frac{1}{1 + t} = \sum_{k<N}(-t)^k + \frac{(-t)^N}{1 + t}$ and integrating, with $\int_0^\infty t^ke^{-xt}dt = k!/x^{k+1}$,
\[
f(x) = \sum_{k=0}^{N-1}\frac{(-1)^kk!}{x^{k+1}} + R_N(x), \qquad |R_N(x)| \leq \int_0^\infty t^Ne^{-xt}dt = \frac{N!}{x^{N+1}}.
\]
So $f(x) \sim \sum_k(-1)^kk!\,x^{-k-1}$, a series that diverges for every $x$. Still, for $x = 10$ the terms decrease until $k \approx 10$, and truncating there gives $f(10)$ with an error below $10!/10^{11} \approx 4 \cdot 10^{-5}$.
\end{example}

\begin{counterexample}\label{special-functions:counterexample-asymptotic-not-unique}
Different functions can have the same asymptotic expansion: $e^{-x} \sim 0$, since $x^Ne^{-x} \to 0$ for every $N$. So $f$ and $f + e^{-x}$ have the same expansion; an asymptotic expansion determines a function only
up to terms smaller than every power.
\end{counterexample}

\begin{theorem}[Watson's lemma]\label{special-functions:theorem-watson}
Let $\varphi$ be locally integrable on $[0, \infty)$ with $|\varphi(t)| \leq Ce^{bt}$ for $t \geq 1$, and suppose that near $0$, $\varphi(t) = \sum_{k<N}a_kt^{k+\alpha-1} + O(t^{N+\alpha-1})$ for every $N$, where $\alpha > 0$.
Then, as $x \to \infty$,
\[
\int_0^\infty e^{-xt}\varphi(t)\,dt \sim \sum_{k \geq 0}a_k\frac{\Gamma(k + \alpha)}{x^{k+\alpha}}.
\]
\end{theorem}

\begin{proof}
Fix $N$, and $\delta \in (0, 1]$ with $|\varphi(t) - \sum_{k<N}a_kt^{k+\alpha-1}| \leq Kt^{N+\alpha-1}$ on $(0, \delta]$. Split the integral at $\delta$. The part over $[\delta, \infty)$ is $O(e^{-x\delta/2})$ for $x \geq 2b + 2$, because
$|e^{-xt}\varphi(t)| \leq e^{-x\delta/2}e^{-(x/2)t}|\varphi(t)|$ there and $\int_\delta^\infty e^{-(x/2)t}|\varphi(t)|dt$ stays bounded. On $(0, \delta]$ replace $\varphi$ by the polynomial part: the error is at most
$K\int_0^\infty t^{N+\alpha-1}e^{-xt}dt = K\Gamma(N + \alpha)x^{-N-\alpha}$, and $\int_0^\delta t^{k+\alpha-1}e^{-xt}dt = \Gamma(k + \alpha)x^{-k-\alpha} - \int_\delta^\infty t^{k+\alpha-1}e^{-xt}dt$, where the last integral is exponentially small. All
exponentially small terms are $O(x^{-N-\alpha})$.
\end{proof}

\begin{theorem}[Laplace's method]\label{special-functions:theorem-laplace}
Let $-\infty \leq a < c < b \leq \infty$, $h$ a real function on $(a, b)$, twice continuously differentiable near $c$, with $h^{\prime}(c) = 0$, $h^{\prime\prime}(c) < 0$, and $\sup_{|t - c| \geq \delta}h(t) < h(c)$ for every
$\delta > 0$. Let $g$ be bounded and continuous at $c$, with $g(c) \neq 0$, and assume $\int_a^be^{\lambda_0h(t)}\,dt < \infty$ for some $\lambda_0$. Then, as $\lambda \to \infty$,
\[
\int_a^bg(t)e^{\lambda h(t)}\,dt \sim g(c)e^{\lambda h(c)}\sqrt{\frac{2\pi}{\lambda|h^{\prime\prime}(c)|}}.
\]
\end{theorem}

\begin{proof}
We may assume $h(c) = 0$. Choose $\delta > 0$ with $h^{\prime\prime} \leq \frac12h^{\prime\prime}(c) = -\kappa < 0$ on $[c - \delta, c + \delta]$, and let $\eta = -\sup_{|t-c| \geq \delta}h > 0$.
\emph{Tails.} For $|t - c| \geq \delta$ and $\lambda \geq \lambda_0$, $e^{\lambda h(t)} = e^{(\lambda - \lambda_0)h(t)}e^{\lambda_0h(t)} \leq e^{-(\lambda - \lambda_0)\eta}e^{\lambda_0h(t)}$, so the tails contribute
$O(e^{-\lambda\eta})$, which is $o(\lambda^{-1/2})$.
\emph{Centre.} Substitute $t = c + v/\sqrt{\lambda}$:
\[
\sqrt{\lambda}\int_{c-\delta}^{c+\delta}g(t)e^{\lambda h(t)}dt = \int_{|v| \leq \delta\sqrt{\lambda}}g(c + v/\sqrt{\lambda})e^{\lambda h(c + v/\sqrt{\lambda})}dv.
\]
By Taylor's theorem $\lambda h(c + v/\sqrt{\lambda}) = \frac12h^{\prime\prime}(\xi)v^2$ for some $\xi$ between $c$ and $t$, which is at most $-\kappa v^2/2$ on the range of integration and tends to $\frac12h^{\prime\prime}(c)v^2$. The integrand
is dominated by $\sup|g|\,e^{-\kappa v^2/2}$, so by dominated convergence the integral tends to $g(c)\int_\RR e^{-|h^{\prime\prime}(c)|v^2/2}dv = g(c)\sqrt{2\pi/|h^{\prime\prime}(c)|}$, the Gaussian integral of
Example~\ref{special-functions:example-gamma-half}.
\end{proof}

\begin{corollary}[Stirling's formula for $\Gamma$]\label{special-functions:corollary-stirling-gamma}
As $x \to \infty$ through real values, $\Gamma(x + 1) \sim \sqrt{2\pi x}\,x^xe^{-x}$.
\end{corollary}

\begin{proof}
Substituting $t = xu$, $\Gamma(x + 1) = \int_0^\infty t^xe^{-t}dt = x^{x+1}\int_0^\infty e^{x(\log u - u)}du$. The function $h(u) = \log u - u$ has its maximum $-1$ at $u = 1$, with $h^{\prime\prime}(1) = -1$, it is bounded away
from $-1$ outside every neighbourhood of $1$, and $\int_0^\infty e^{h(u)}du = 1$. Theorem~\ref{special-functions:theorem-laplace} with $g = 1$ gives $\int_0^\infty e^{xh(u)}du \sim e^{-x}\sqrt{2\pi/x}$.
\end{proof}

For integers this is Theorem~\ref{real-analysis:theorem-stirling}. The method of Laplace has two relatives. When the exponent is purely imaginary, $\int g(t)e^{i\lambda h(t)}dt$, the main contributions come from the
stationary points $h^{\prime}(c) = 0$, by cancellation elsewhere, and each nondegenerate one contributes $g(c)e^{i\lambda h(c) \pm i\pi/4}\sqrt{2\pi/(\lambda|h^{\prime\prime}(c)|)}$ (\emph{stationary phase}). When $g$ and $h$
are holomorphic, one deforms the contour through a saddle point of $h$ along the direction of steepest descent of $\operatorname{Re}h$, where $\operatorname{Im}h$ is constant, which reduces to Laplace's method
(\emph{saddle point method}) \cite{olver-asymptotics}.

\begin{remark}\label{special-functions:remark-stationary-phase}
\emph{Stationary phase.} If $h$ is smooth on $[a, b]$ with a single stationary point $c \in (a, b)$, $h^{\prime\prime}(c) \neq 0$, and $g$ is smooth with compact support in $(a, b)$, then
$\int_a^bg(t)e^{i\lambda h(t)}dt = g(c)e^{i\lambda h(c) + i\frac\pi4\operatorname{sgn}h^{\prime\prime}(c)}\sqrt{\frac{2\pi}{\lambda|h^{\prime\prime}(c)|}} + O(\lambda^{-1})$. Away from $c$ one integrates by parts, writing
$e^{i\lambda h} = (i\lambda h^{\prime})^{-1}\frac{d}{dt}e^{i\lambda h}$, which gains a factor $\lambda^{-1}$ each time; near $c$ one changes variables to make $h$ exactly quadratic and uses the Fresnel integral
$\int_\RR e^{i\lambda v^2/2}dv = e^{i\pi/4}\sqrt{2\pi/\lambda}$ \cite{olver-asymptotics}. It gives the asymptotics of Bessel functions (Remark~\ref{special-functions:remark-bessel-asymptotics}).
\end{remark}

The other basic tool compares sums with integrals.

\begin{definition}\label{special-functions:definition-bernoulli-polynomials}
The \emph{Bernoulli polynomials} $B_n(x)$ are defined by $\frac{te^{xt}}{e^t - 1} = \sum_{n \geq 0}B_n(x)\frac{t^n}{n!}$. So $B_n(0) = B_n$ are the Bernoulli numbers of
Definition~\ref{generating-functions:definition-bernoulli}, $B_0(x) = 1$, $B_1(x) = x - \frac12$, $B_2(x) = x^2 - x + \frac16$. The \emph{periodic Bernoulli functions} are $P_n(x) = B_n(x - \lfloor x\rfloor)$.
\end{definition}

\begin{lemma}\label{special-functions:lemma-bernoulli-polynomials}
$B_n^{\prime}(x) = nB_{n-1}(x)$ for $n \geq 1$; $B_n(1) = B_n(0)$ for $n \geq 2$; and $B_n = 0$ for odd $n \geq 3$. Hence $P_n$ is continuous for $n \geq 2$, and $P_n^{\prime} = nP_{n-1}$ off the integers for $n \geq 1$.
\end{lemma}

\begin{proof}
Differentiating the generating function in $x$ gives $\sum B_n^{\prime}(x)t^n/n! = t\sum B_n(x)t^n/n!$, which is the first claim. Next $\frac{te^t}{e^t - 1} - \frac{t}{e^t - 1} = t$, so $B_n(1) - B_n(0)$ is $0$ for $n \neq 1$. Finally
$\frac{t}{e^t - 1} + \frac t2 = \frac t2\frac{e^t + 1}{e^t - 1}$ is even in $t$, so the odd coefficients other than $B_1$ vanish.
\end{proof}

\begin{theorem}[Euler--Maclaurin summation]\label{special-functions:theorem-euler-maclaurin}
Let $a < b$ be integers, $m \geq 1$ and $f \in C^m[a, b]$. Then
\[
\sum_{a < n \leq b}f(n) = \int_a^bf(x)\,dx + \sum_{k=1}^m\frac{(-1)^kB_k}{k!}\bigl(f^{(k-1)}(b) - f^{(k-1)}(a)\bigr) - \frac{(-1)^m}{m!}\int_a^bP_m(x)f^{(m)}(x)\,dx.
\]
\end{theorem}

\begin{proof}
\emph{The case $m = 1$.} On $[n, n + 1]$, $P_1(x) = x - n - \frac12$, and integration by parts gives $\int_n^{n+1}P_1f^{\prime} = \frac12(f(n) + f(n + 1)) - \int_n^{n+1}f$. Summing over $a \leq n < b$,
$\int_a^bP_1f^{\prime} = \sum_{a<n\leq b}f(n) - \frac12(f(b) - f(a)) - \int_a^bf$, which is the formula, as $-B_1 = \frac12$.
\emph{Induction.} By Lemma~\ref{special-functions:lemma-bernoulli-polynomials}, $P_{m+1}$ is continuous with derivative $(m + 1)P_m$ and $P_{m+1}(n) = B_{m+1}$ at integers, so
$\int_a^bP_mf^{(m)} = \frac{B_{m+1}}{m + 1}\bigl(f^{(m)}(b) - f^{(m)}(a)\bigr) - \frac{1}{m + 1}\int_a^bP_{m+1}f^{(m+1)}$. Multiplying by $-(-1)^m/m!$ turns the remainder for $m$ into the term $k = m + 1$ plus the
remainder for $m + 1$.
\end{proof}

\begin{example}\label{special-functions:example-harmonic-expansion}
For $f(x) = 1/x$ on $[1, N]$ with $m = 2$, adding $f(1) = 1$ and using $-B_1 = \frac12$, $\frac{B_2}{2} = \frac{1}{12}$,
\[
H_N = \log N + \frac12 + \frac{1}{2N} + \frac1{12}\Bigl(1 - \frac{1}{N^2}\Bigr) - \int_1^N\frac{P_2(x)}{x^3}dx.
\] Letting $N \to \infty$ identifies $\gamma = \frac12 + \frac1{12} - \int_1^\infty P_2x^{-3}dx$, and
\[
H_N = \log N + \gamma + \frac{1}{2N} - \frac{1}{12N^2} + \int_N^\infty\frac{P_2(x)}{x^3}dx = \log N + \gamma + \frac{1}{2N} - \frac{1}{12N^2} + O(N^{-4}),
\]
the last step by integrating by parts twice, using $P_2 = \frac13P_3^{\prime}$, $P_3 = \frac14P_4^{\prime}$ and $B_3 = 0$. Continuing gives $H_N \sim \log N + \gamma + \frac{1}{2N} - \sum_{k \geq 1}\frac{B_{2k}}{2kN^{2k}}$, divergent like the series of Example~\ref{special-functions:example-stieltjes}.
In the same way $\log n! = (n + \frac12)\log n - n + \frac12\log 2\pi + \sum_{k=1}^{K}\frac{B_{2k}}{2k(2k - 1)n^{2k-1}} + O(n^{-2K-1})$, the \emph{Stirling series}.
\end{example}

\section{Orthogonal polynomials}\label{special-functions:section-orthogonal-polynomials}

\begin{definition}\label{special-functions:definition-orthogonal-polynomials}
Let $(a, b)$ be an interval ($a$, $b$ possibly infinite) and $w \geq 0$ a \emph{weight}: a measurable function, positive on a set of positive measure, with $\int_a^b|x|^nw(x)\,dx < \infty$ for all $n$. The inner product
$\langle f, g\rangle = \int_a^bf(x)g(x)w(x)\,dx$ is positive definite on real polynomials. The \emph{monic orthogonal polynomials} $p_0, p_1, \ldots$ are obtained by Gram--Schmidt
(Theorem~\ref{linear-algebra:theorem-gram-schmidt}) from $1, x, x^2, \ldots$ without normalising: $p_n$ is monic of degree $n$ and orthogonal to all polynomials of degree $< n$.
\end{definition}

\begin{notation}\label{special-functions:notation-orthogonal}
$\langle\ ,\ \rangle$ and $\|\ \|$ are the inner product and norm of Definition~\ref{special-functions:definition-orthogonal-polynomials}; $\mathcal{P}_n$ is the space of real polynomials of degree at most $n$. The classical
families below are normalised differently (not monic), following tradition.
\end{notation}

\begin{theorem}[Three-term recurrence]\label{special-functions:theorem-three-term}
There are real numbers $\alpha_n$ and $\beta_n > 0$ with $p_{n+1}(x) = (x - \alpha_n)p_n(x) - \beta_np_{n-1}(x)$ for $n \geq 1$, namely $\alpha_n = \langle xp_n, p_n\rangle/\|p_n\|^2$ and $\beta_n = \|p_n\|^2/\|p_{n-1}\|^2$;
also $p_1 = x - \alpha_0$.
\end{theorem}

\begin{proof}
$xp_n - p_{n+1}$ has degree at most $n$, so it is $\sum_{k \leq n}c_kp_k$ with $c_k\|p_k\|^2 = \langle xp_n, p_k\rangle = \langle p_n, xp_k\rangle$, using that multiplication by $x$ is symmetric. For $k \leq n - 2$, $xp_k$ has
degree $< n$, so $c_k = 0$. For $k = n$ this gives $\alpha_n$; for $k = n - 1$, $xp_{n-1} = p_n + (\text{lower degree})$, so $\langle p_n, xp_{n-1}\rangle = \|p_n\|^2$ and $c_{n-1} = \beta_n$.
\end{proof}

\begin{theorem}\label{special-functions:theorem-zeros}
$p_n$ has $n$ distinct real zeros, all in $(a, b)$.
\end{theorem}

\begin{proof}
Let $x_1 < \cdots < x_m$ be the points of $(a, b)$ where $p_n$ changes sign, and $q = \prod_i(x - x_i)$. Then $p_nq$ does not change sign on $(a, b)$ and is not identically zero, so $\langle p_n, q\rangle \neq 0$. As $p_n$ is
orthogonal to $\mathcal{P}_{n-1}$, $\deg q = m \geq n$; and $m \leq n$ since $p_n$ has at most $n$ zeros. So $p_n$ has $n$ sign changes, hence $n$ simple zeros, in $(a, b)$.
\end{proof}

\begin{theorem}[Gauss quadrature]\label{special-functions:theorem-gauss-quadrature}
Let $x_1, \ldots, x_n$ be the zeros of $p_n$. There are unique weights $\lambda_1, \ldots, \lambda_n$ such that $\int_a^bq(x)w(x)\,dx = \sum_i\lambda_iq(x_i)$ for all $q \in \mathcal{P}_{2n-1}$. The weights are positive, and
no choice of $n$ nodes and weights is exact on $\mathcal{P}_{2n}$.
\end{theorem}

\begin{proof}
Let $\ell_i$ be the Lagrange polynomials of degree $n - 1$ with $\ell_i(x_j) = \delta_{ij}$, and $\lambda_i = \int\ell_iw$; then the rule is exact on $\mathcal{P}_{n-1}$, and these are the only weights with this property (apply the
rule to $\ell_j$). For $q \in \mathcal{P}_{2n-1}$ divide: $q = sp_n + r$ with $s, r \in \mathcal{P}_{n-1}$. Then $\int qw = \langle s, p_n\rangle + \int rw = \int rw = \sum_i\lambda_ir(x_i) = \sum_i\lambda_iq(x_i)$, since
$p_n(x_i) = 0$. Applied to $q = \ell_j^2 \in \mathcal{P}_{2n-2}$ this gives $\lambda_j = \int\ell_j^2w > 0$. For any nodes $y_1, \ldots, y_n$, the polynomial $\prod(x - y_i)^2 \in \mathcal{P}_{2n}$ has positive integral and vanishes at
the nodes.
\end{proof}

\begin{example}\label{special-functions:example-classical-polynomials}
The classical families, each orthogonal for the weight shown:
\begin{enumerate}
\item \emph{Legendre} polynomials, on $(-1, 1)$ with $w = 1$: $P_n(x) = \frac{1}{2^nn!}\frac{d^n}{dx^n}(x^2 - 1)^n$ (Rodrigues' formula), $P_0 = 1$, $P_1 = x$, $P_2 = \frac12(3x^2 - 1)$, with
$(n + 1)P_{n+1} = (2n + 1)xP_n - nP_{n-1}$, $\int_{-1}^1P_n^2 = \frac{2}{2n + 1}$, $P_n(1) = 1$, and generating function $\sum_nP_n(x)t^n = (1 - 2xt + t^2)^{-1/2}$. They were found by Legendre in the expansion of the
Newtonian potential $1/|x - y|$; they reappear in Functional Analysis as an orthonormal basis of $L^2(-1, 1)$ after normalisation.
\item \emph{Chebyshev} polynomials, on $(-1, 1)$ with $w = (1 - x^2)^{-1/2}$: $T_n(\cos\theta) = \cos n\theta$, $T_{n+1} = 2xT_n - T_{n-1}$, zeros $\cos\frac{(2k - 1)\pi}{2n}$.
\item \emph{Hermite} polynomials, on $\RR$ with $w = e^{-x^2}$: $H_n(x) = (-1)^ne^{x^2}\frac{d^n}{dx^n}e^{-x^2}$, $H_{n+1} = 2xH_n - 2nH_{n-1}$, $\int H_n^2e^{-x^2} = 2^nn!\sqrt{\pi}$, $\sum_nH_n(x)\frac{t^n}{n!} = e^{2xt - t^2}$.
\item \emph{Laguerre} polynomials, on $(0, \infty)$ with $w = e^{-x}$: $L_n(x) = \frac{e^x}{n!}\frac{d^n}{dx^n}(e^{-x}x^n)$, $\sum_nL_n(x)t^n = \frac{1}{1 - t}e^{-xt/(1-t)}$.
\item \emph{Jacobi} polynomials $P_n^{(\alpha, \beta)}$, on $(-1, 1)$ with $w = (1 - x)^\alpha(1 + x)^\beta$, $\alpha, \beta > -1$, contain the Legendre ($\alpha = \beta = 0$) and Chebyshev ($\alpha = \beta = -\frac12$) cases \cite{szego-orthogonal}.
\end{enumerate}
\end{example}

\begin{proposition}\label{special-functions:proposition-legendre-hermite}
\begin{enumerate}
\item The polynomials $P_n$ given by Rodrigues' formula are orthogonal on $(-1, 1)$ with $w = 1$, $P_n(1) = 1$, and $\int_{-1}^1P_n^2 = \frac{2}{2n + 1}$.
\item The Hermite polynomials satisfy $\sum_nH_n(x)t^n/n! = e^{2xt - t^2}$ and $\int_\RR H_mH_ne^{-x^2}dx = \delta_{mn}2^nn!\sqrt{\pi}$.
\item $T_n(x) = \cos(n\arccos x)$ is a polynomial of degree $n$ with leading coefficient $2^{n-1}$ ($n \geq 1$), and $\int_{-1}^1T_mT_n(1 - x^2)^{-1/2}dx = 0$ for $m \neq n$.
\end{enumerate}
\end{proposition}

\begin{proof}
(1) $u = (x^2 - 1)^n$ and its first $n - 1$ derivatives vanish at $\pm 1$. For $q$ of degree $< n$, integrating by parts $n$ times gives $\int_{-1}^1q\,u^{(n)} = (-1)^n\int_{-1}^1q^{(n)}u = 0$. By Leibniz' rule applied to
$u = (x - 1)^n(x + 1)^n$, $u^{(n)}(1) = n!\,2^n$, so $P_n(1) = 1$. For the norm, $P_n$ has leading coefficient $\frac{(2n)!}{2^n(n!)^2}$, so $\int P_n^2 = \frac{(2n)!}{2^n(n!)^2}\int x^nP_n = \frac{(2n)!}{2^{2n}(n!)^3}(-1)^nn!\int_{-1}^1(x^2 - 1)^n$,
and, substituting $x^2 = u$ and using Proposition~\ref{special-functions:proposition-beta}, $\int_{-1}^1(1 - x^2)^ndx = B(\frac12, n + 1) = \frac{\Gamma(\frac12)\,n!}{\Gamma(n + \frac32)}$. From $\Gamma(\frac12) = \sqrt\pi$
and the functional equation, $\Gamma(n + \frac32) = \frac{(2n + 2)!}{4^{n+1}(n + 1)!}\sqrt\pi$, so $\int_{-1}^1(1 - x^2)^ndx = \frac{2^{2n+1}(n!)^2}{(2n + 1)!}$, and the product is $\frac{2}{2n + 1}$.
(2) $e^{2xt - t^2} = e^{x^2}e^{-(x - t)^2}$, whose $n$-th $t$-derivative at $t = 0$ is $e^{x^2}(-1)^n\frac{d^n}{dx^n}e^{-x^2} = H_n(x)$. Multiplying two generating functions,
$\int e^{2xt - t^2}e^{2xs - s^2}e^{-x^2}dx = e^{2st}\int e^{-(x - s - t)^2}dx = \sqrt{\pi}e^{2st}$, and comparing coefficients of $t^ms^n$ gives the orthogonality relations.
(3) $\cos(n + 1)\theta + \cos(n - 1)\theta = 2\cos\theta\cos n\theta$ gives $T_{n+1} = 2xT_n - T_{n-1}$, with $T_0 = 1$, $T_1 = x$, and induction gives the degree and leading coefficient. With $x = \cos\theta$,
$\int_{-1}^1T_mT_n(1 - x^2)^{-1/2}dx = \int_0^\pi\cos m\theta\cos n\theta\,d\theta = 0$ for $m \neq n$.
\end{proof}

\begin{proposition}[Chebyshev]\label{special-functions:proposition-chebyshev-minimal}
Among monic real polynomials of degree $n \geq 1$, $2^{1-n}T_n$ has the smallest maximum of $|q|$ on $[-1, 1]$, namely $2^{1-n}$.
\end{proposition}

\begin{proof}
$2^{1-n}T_n$ is monic with $\max|2^{1-n}T_n| = 2^{1-n}$, attained with alternating signs at the $n + 1$ points $\cos(k\pi/n)$, $k = 0, \ldots, n$. If $q$ is monic with $\max|q| < 2^{1-n}$, then $r = 2^{1-n}T_n - q$
has degree $< n$ and alternates in sign at these $n + 1$ points, so it has $n$ zeros, hence $r = 0$, a contradiction.
\end{proof}

This is why Chebyshev nodes are good interpolation nodes: the interpolation error of a smooth function at the nodes $x_i$ is $\frac{f^{(n)}(\xi)}{n!}\prod(x - x_i)$, and $\prod(x - x_i)$ is smallest for the zeros of $T_n$.

\begin{counterexample}\label{special-functions:counterexample-runge}
Equally spaced nodes are bad: interpolating $f(x) = \frac{1}{1 + 25x^2}$ at $n$ equally spaced points of $[-1, 1]$, the maximal error tends to infinity as $n \to \infty$ (\emph{Runge's phenomenon}), because $f$ has poles at
$\pm i/5$, close to the interval. At Chebyshev nodes the interpolants converge uniformly for every function that is analytic near $[-1, 1]$ \cite{trefethen-approximation}.
\end{counterexample}

\section{Linear differential equations in the complex domain}\label{special-functions:section-fuchsian}

Most special functions solve second-order linear differential equations with holomorphic coefficients. We study such equations near ordinary and singular points.

\begin{theorem}\label{special-functions:theorem-holomorphic-solutions}
Let $A(z)$ be an $n \times n$ matrix of holomorphic functions on the disc $|z| < R$ and $w_0 \in \CC^n$. The system $w^{\prime} = A(z)w$ has a unique holomorphic solution on $|z| < R$ with $w(0) = w_0$.
\end{theorem}

\begin{proof}
\emph{Formal solution.} Write $A(z) = \sum_jA_jz^j$ and $w = \sum_kw_kz^k$. The equation is equivalent to $(k + 1)w_{k+1} = \sum_{j=0}^kA_jw_{k-j}$, which determines the $w_k$ uniquely from $w_0$; so a holomorphic
solution is unique if it exists.
\emph{Majorant.} Fix $r < R$ and $M$ with $\|A_j\| \leq Mr^{-j}$ (the Cauchy estimates of Theorem~\ref{complex-analysis:theorem-cauchy-formula}, applied entrywise). Let $c(z) = |w_0|(1 - z/r)^{-Mr} = \sum_kc_kz^k$, which solves
$c^{\prime} = M(1 - z/r)^{-1}c$, that is, $(k + 1)c_{k+1} = \sum_jMr^{-j}c_{k-j}$, with $c_k \geq 0$. By induction $|w_k| \leq c_k$: indeed
$(k + 1)|w_{k+1}| \leq \sum_jMr^{-j}|w_{k-j}| \leq \sum_jMr^{-j}c_{k-j} = (k + 1)c_{k+1}$. The series for $c$ has radius $r$, so $\sum_kw_kz^k$ converges on $|z| < r$, and on $|z| < R$ since $r < R$ was arbitrary.
\end{proof}

A linear equation $w^{(n)} + a_{n-1}(z)w^{(n-1)} + \cdots + a_0(z)w = 0$ is the system for $(w, w^{\prime}, \ldots, w^{(n-1)})$, so it has an $n$-dimensional space of holomorphic solutions on every disc where the coefficients are
holomorphic. On a domain $U$ that is not simply connected, solutions can be continued along paths, but continuation around a loop may not return to the same solution.

\begin{definition}\label{special-functions:definition-monodromy}
Let the coefficients be holomorphic on a domain $U$, $z_0 \in U$, and $V$ the space of solutions near $z_0$. Analytic continuation along a loop $\gamma$ at $z_0$ defines a linear map $M_\gamma \colon V \to V$, depending only on the
homotopy class of $\gamma$; the resulting homomorphism $\pi_1(U, z_0) \to \GL(V)$ (up to the convention on composition) is the \emph{monodromy representation} of the equation.
\end{definition}

That $M_\gamma$ depends only on the homotopy class is the monodromy theorem: continuation along a homotopy is possible because of Theorem~\ref{special-functions:theorem-holomorphic-solutions} applied on small discs, and the
result varies locally constantly, as in the proof of homotopy lifting for coverings (Proposition~\ref{homotopy:proposition-homotopy-lifting}).

\begin{example}\label{special-functions:example-monodromy}
On $U = \CC \setminus \{0\}$, $zw^{\prime} = \alpha w$ has the solution $z^\alpha = e^{\alpha\log z}$, and continuation once around $0$ multiplies it by $e^{2\pi i\alpha}$. The equation $(zw^{\prime})^{\prime} = 0$, that is
$w^{\prime\prime} + \frac1zw^{\prime} = 0$, has solutions $1$ and $\log z$, and the monodromy is $\log z \mapsto \log z + 2\pi i$, the unipotent matrix $\begin{pmatrix} 1 & 2\pi i \\ 0 & 1 \end{pmatrix}$.
\end{example}

\begin{definition}\label{special-functions:definition-regular-singular}
For $w^{\prime\prime} + p(z)w^{\prime} + q(z)w = 0$ with $p$, $q$ meromorphic at $0$, the point $0$ is a \emph{regular singular point} if $zp(z)$ and $z^2q(z)$ are holomorphic at $0$, with values $p_0$, $q_0$ there. The
\emph{indicial polynomial} is $F(\rho) = \rho(\rho - 1) + p_0\rho + q_0$, and its roots $\rho_1$, $\rho_2$ are the \emph{exponents} at $0$. The point $\infty$ is regular singular if $0$ is so after the substitution $z = 1/u$.
An equation whose singular points in $\CC \cup \{\infty\}$ are all regular is \emph{Fuchsian}.
\end{definition}

\begin{theorem}[Frobenius]\label{special-functions:theorem-frobenius}
Let $0$ be a regular singular point, with $zp$ and $z^2q$ holomorphic on $|z| < R$, and exponents numbered so that $\operatorname{Re}\rho_1 \geq \operatorname{Re}\rho_2$.
\begin{enumerate}
\item There is a solution $w_1 = z^{\rho_1}\sum_{k \geq 0}c_kz^k$ with $c_0 = 1$, the series converging for $|z| < R$.
\item If $\rho_1 - \rho_2 \notin \ZZ$, there is a second solution $w_2 = z^{\rho_2}\sum_kd_kz^k$, $d_0 = 1$, of the same kind.
\item If $\rho_1 - \rho_2 \in \ZZ_{\geq 0}$, there is a second solution $w_2 = cw_1\log z + z^{\rho_2}\sum_kd_kz^k$ for a constant $c$, which is nonzero if $\rho_1 = \rho_2$.
\end{enumerate}
\end{theorem}

\begin{proof}[Sketch of proof]
Write $zp = \sum p_jz^j$ and $z^2q = \sum q_jz^j$, and multiply the equation by $z^2$. Substituting $w = z^\rho\sum c_kz^k$ gives $F(\rho + k)c_k = -\sum_{j=1}^{k}\bigl((\rho + k - j)p_j + q_j\bigr)c_{k-j}$. For
$\rho = \rho_1$, $F(\rho_1 + k) = k(k + \rho_1 - \rho_2) \neq 0$ for $k \geq 1$, because $\operatorname{Re}(\rho_1 - \rho_2) \geq 0$; so the $c_k$ are determined. Since $|F(\rho_1 + k)| \geq k(k - |\rho_1 - \rho_2|)$ grows like $k^2$
while the right side involves $k$ linearly, a majorant argument as in Theorem~\ref{special-functions:theorem-holomorphic-solutions} shows convergence on $|z| < R$. The same works for $\rho_2$ when $F(\rho_2 + k) \neq 0$ for all
$k \geq 1$, that is, when $\rho_1 - \rho_2 \notin \ZZ_{>0}$; in case (3), reduction of order, $w_2 = w_1\int w_1^{-2}e^{-\int p}$, produces the logarithm. The omitted verifications are the majorant estimates
\cite{whittaker-watson}.
\end{proof}

\begin{example}\label{special-functions:example-frobenius}
\begin{enumerate}
\item The Euler equation $z^2w^{\prime\prime} + p_0zw^{\prime} + q_0w = 0$ has the solutions $z^{\rho_1}$, $z^{\rho_2}$, and $z^\rho$, $z^\rho\log z$ when $\rho_1 = \rho_2 = \rho$.
\item Bessel's equation $z^2w^{\prime\prime} + zw^{\prime} + (z^2 - \nu^2)w = 0$ has exponents $\pm\nu$ at $0$, and the Frobenius series for $\rho = \nu$ is the Bessel function of Definition~\ref{special-functions:definition-bessel}.
\item Legendre's equation $(1 - z^2)w^{\prime\prime} - 2zw^{\prime} + \lambda w = 0$ has regular singular points at $\pm 1$ and $\infty$, with exponents $0, 0$ at $\pm 1$. A solution holomorphic at both $1$ and $-1$ exists only if
$\lambda = n(n + 1)$, and then it is the Legendre polynomial $P_n$; this is why Legendre polynomials appear when one separates variables in the Laplace equation in spherical coordinates.
\end{enumerate}
\end{example}

\begin{counterexample}\label{special-functions:counterexample-irregular}
At an irregular singular point the theorem fails. The equation $z^2w^{\prime} + w = z$ has the formal solution $\sum_{k \geq 0}(-1)^kk!\,z^{k+1}$, with radius of convergence $0$; an actual solution near $0$ is
$w(z) = \int_0^\infty\frac{e^{-t/z}}{1 + t}dt$ for $z > 0$, which has this series as its asymptotic expansion (Example~\ref{special-functions:example-stieltjes} with $x = 1/z$). And $z^2w^{\prime} = w$ has the solution $e^{-1/z}$,
with an essential singularity: no solution of the form $z^\rho\sum c_kz^k$.
\end{counterexample}

\begin{remark}\label{special-functions:remark-riemann-p}
For a second-order Fuchsian equation with singular points $z_1, \ldots, z_m$ (including $\infty$), the sum of all the exponents is $m - 2$ (\emph{Fuchs' relation}). With three singular points, which a M\"obius transformation
moves to $0, 1, \infty$, Riemann showed that the equation is determined by its six exponents, and after the substitution $w \mapsto z^\lambda(1 - z)^\mu w$ it becomes the hypergeometric equation of the next section
\cite{whittaker-watson}. Hilbert's 21st problem asks, conversely, for a Fuchsian system with prescribed singular points and monodromy; this is the Riemann--Hilbert correspondence.
\end{remark}

\section{The hypergeometric function}\label{special-functions:section-hypergeometric}

\begin{definition}\label{special-functions:definition-hypergeometric}
The \emph{Pochhammer symbol} is $(a)_0 = 1$, $(a)_n = a(a + 1)\cdots(a + n - 1)$, so $(a)_n = \Gamma(a + n)/\Gamma(a)$ when $a \notin \ZZ_{\leq 0}$. For $c \notin \ZZ_{\leq 0}$, the \emph{hypergeometric function} is
\[
F(a, b; c; z) = {}_2F_1(a, b; c; z) = \sum_{n \geq 0}\frac{(a)_n(b)_n}{(c)_n\,n!}z^n, \qquad |z| < 1.
\]
\end{definition}

The ratio of consecutive coefficients is $\frac{(a + n)(b + n)}{(c + n)(n + 1)} \to 1$, so the radius of convergence is $1$ unless $a$ or $b$ is a nonpositive integer, in which case $F$ is a polynomial.

\begin{example}\label{special-functions:example-hypergeometric}
\begin{enumerate}
\item $F(a, b; b; z) = \sum\frac{(a)_n}{n!}z^n = (1 - z)^{-a}$, the binomial series.
\item $zF(1, 1; 2; -z) = \sum\frac{(-1)^nz^{n+1}}{n + 1} = \log(1 + z)$, and $zF(\frac12, \frac12; \frac32; z^2) = \arcsin z$.
\item $P_n(x) = F(-n, n + 1; 1; \frac{1 - x}{2})$ for the Legendre polynomials (Exercise~\ref{special-functions:exercise-legendre-hypergeometric}).
\item $K(k) = \frac\pi2F(\frac12, \frac12; 1; k^2)$ for the complete elliptic integral of Definition~\ref{special-functions:definition-elliptic-integrals}, by expanding $(1 - k^2\sin^2\theta)^{-1/2}$ and integrating termwise.
\end{enumerate}
\end{example}

\begin{proposition}\label{special-functions:proposition-hypergeometric-equation}
$F(a, b; c; z)$ satisfies the \emph{hypergeometric equation}
\[
z(1 - z)w^{\prime\prime} + \bigl(c - (a + b + 1)z\bigr)w^{\prime} - abw = 0,
\]
a Fuchsian equation with singular points $0, 1, \infty$ and exponents $\{0, 1 - c\}$, $\{0, c - a - b\}$, $\{a, b\}$. If $c \notin \ZZ$, a second solution near $0$ is $z^{1-c}F(a - c + 1, b - c + 1; 2 - c; z)$.
\end{proposition}

\begin{proof}
With $\theta = z\frac{d}{dz}$, which acts on $z^n$ as multiplication by $n$, the equation multiplied by $z$ reads $\theta(\theta + c - 1)w = z(\theta + a)(\theta + b)w$. On $w = \sum A_nz^n$ this is
$n(n + c - 1)A_n = (n - 1 + a)(n - 1 + b)A_{n-1}$, which is the recursion satisfied by $A_n = \frac{(a)_n(b)_n}{(c)_nn!}$. The exponents are read off from the leading terms at each point (at $\infty$, substitute $z = 1/u$),
and they add up to $1$, in accordance with Fuchs' relation. Substituting $w = z^{1-c}v$ turns the equation into the hypergeometric equation with parameters $(a - c + 1, b - c + 1, 2 - c)$.
\end{proof}

\begin{theorem}[Euler's integral]\label{special-functions:theorem-euler-integral}
If $\operatorname{Re}c > \operatorname{Re}b > 0$ and $|z| < 1$, then
\[
F(a, b; c; z) = \frac{\Gamma(c)}{\Gamma(b)\Gamma(c - b)}\int_0^1t^{b-1}(1 - t)^{c-b-1}(1 - zt)^{-a}\,dt.
\]
The right side is holomorphic on $\CC \setminus [1, \infty)$, which continues $F$ there.
\end{theorem}

\begin{proof}
For $|u| < 1$, $(1 - u)^{-a} = \sum_n\frac{(a)_n}{n!}u^n$ (Example~\ref{special-functions:example-hypergeometric}(1): both sides are holomorphic, equal to $1$ at $0$, and satisfy $(1 - u)f^{\prime} = af$, so $(1 - u)^af$ has derivative $0$).
For $|zt| \leq |z| < 1$ the series converges uniformly in $t \in [0, 1]$, so we may integrate termwise, and by Proposition~\ref{special-functions:proposition-beta}
\[
\int_0^1t^{b+n-1}(1 - t)^{c-b-1}dt = \frac{\Gamma(b + n)\Gamma(c - b)}{\Gamma(c + n)} = \frac{\Gamma(b)\Gamma(c - b)}{\Gamma(c)}\cdot\frac{(b)_n}{(c)_n}.
\]
Summing gives the formula. For $z \notin [1, \infty)$, $1 - zt \neq 0$ for $t \in [0, 1]$, and Lemma~\ref{special-functions:lemma-holomorphic-integrals} gives holomorphy.
\end{proof}

\begin{theorem}[Gauss]\label{special-functions:theorem-gauss-summation}
If $\operatorname{Re}c > \operatorname{Re}b > 0$ and $\operatorname{Re}(c - a - b) > 0$, the series $F(a, b; c; 1)$ converges absolutely and
\[
F(a, b; c; 1) = \frac{\Gamma(c)\Gamma(c - a - b)}{\Gamma(c - a)\Gamma(c - b)}.
\]
\end{theorem}

\begin{proof}
\emph{Convergence.} The partial products of the Weierstrass product (Theorem~\ref{special-functions:theorem-weierstrass-gamma}) are $\frac{(s)_{n+1}}{n!\,n^s}e^{s(\gamma - H_n + \log n)}$, so
$\frac{(s)_{n+1}}{n!\,n^s} \to \frac{1}{\Gamma(s)}$ for every $s \in \CC$ (Euler's limit formula without the restriction of Theorem~\ref{special-functions:theorem-euler-limit}); multiplying by $\frac{n}{s + n}$,
$(s)_n/((n - 1)!\,n^s) \to 1/\Gamma(s)$. So the $n$-th coefficient is $O(n^{\operatorname{Re}(a + b - c) - 1})$ (it vanishes for large $n$ if $a$ or $b$ is a nonpositive integer), the series converges absolutely on $|z| \leq 1$, and $F(a, b; c; z) \to F(a, b; c; 1)$
as $z \to 1^-$ by uniform convergence on $[0, 1]$.
\emph{Value.} For $0 < z < 1$ and $t \in [0, 1)$, $1 - zt \geq 1 - t > 0$, so $|(1 - zt)^{-a}| = (1 - zt)^{-\operatorname{Re}a} \leq 1 + (1 - t)^{-\operatorname{Re}a}$. This makes the integrand of Theorem~\ref{special-functions:theorem-euler-integral}
dominated by an integrable function, because $\operatorname{Re}(c - b) > 0$ and $\operatorname{Re}(c - b - a) > 0$. By dominated convergence, as $z \to 1^-$ the integral tends to
$\int_0^1t^{b-1}(1 - t)^{c-a-b-1}dt = \frac{\Gamma(b)\Gamma(c - a - b)}{\Gamma(c - a)}$, and the formula follows.
\end{proof}

\begin{corollary}[Chu--Vandermonde]\label{special-functions:corollary-chu-vandermonde}
For integers $n \geq 0$ and $c \notin \ZZ_{\leq 0}$, $F(-n, b; c; 1) = \frac{(c - b)_n}{(c)_n}$.
\end{corollary}

\begin{proof}
Both sides are rational functions of $b$ and $c$. For $\operatorname{Re}c > \operatorname{Re}b > 0$, Theorem~\ref{special-functions:theorem-gauss-summation} applies with $a = -n$ (then $\operatorname{Re}(c - a - b) > n \geq 0$) and gives
$\frac{\Gamma(c)\Gamma(c - b + n)}{\Gamma(c + n)\Gamma(c - b)} = \frac{(c - b)_n}{(c)_n}$; two rational functions that agree on an open set agree.
\end{proof}

\begin{remark}\label{special-functions:remark-kummer-confluent}
Kummer found $24$ solutions of the hypergeometric equation of the form $z^\lambda(1 - z)^\mu F(\cdots; \cdots; \phi(z))$, with $\phi$ one of the six M\"obius maps permuting $0, 1, \infty$; for example
$F(a, b; c; z) = (1 - z)^{-a}F(a, c - b; c; \frac{z}{z - 1})$ (\emph{Pfaff's transformation}), which follows from Euler's integral by $t \mapsto 1 - t$. Letting $b \to \infty$ in $F(a, b; c; z/b)$ gives the \emph{confluent
hypergeometric function} $M(a; c; z) = {}_1F_1(a; c; z) = \sum_n\frac{(a)_n}{(c)_nn!}z^n$, an entire function solving Kummer's equation $zw^{\prime\prime} + (c - z)w^{\prime} - aw = 0$, with an irregular singular point at $\infty$;
it contains $e^z = M(a; a; z)$, the error function and the Hermite and Laguerre polynomials \cite{andrews-askey-roy}.
\end{remark}

\section{Bessel functions and spherical harmonics}\label{special-functions:section-bessel}

\begin{definition}\label{special-functions:definition-bessel}
For $\nu \in \CC$, $\nu \notin \ZZ_{<0}$, the \emph{Bessel function of the first kind} is
\[
J_\nu(z) = \sum_{k \geq 0}\frac{(-1)^k}{k!\,\Gamma(k + \nu + 1)}\Bigl(\frac z2\Bigr)^{2k+\nu},
\]
where $(z/2)^\nu$ uses the principal branch; for $\nu = -n$ a negative integer the same formula, with $1/\Gamma = 0$ at the poles, gives $J_{-n} = (-1)^nJ_n$.
\end{definition}

The series $z^{-\nu}J_\nu(z)$ converges for all $z$ (the coefficients decay like $1/(k!)^2$), and the computation of Example~\ref{special-functions:example-frobenius}(2) shows that $J_\nu$ solves Bessel's equation. For
$\nu \notin \ZZ$, $J_\nu$ and $J_{-\nu}$ are independent solutions; for integers one needs a second solution with a logarithm, the Bessel function of the second kind $Y_n$ (Theorem~\ref{special-functions:theorem-frobenius}(3)).

\begin{theorem}[Generating function]\label{special-functions:theorem-bessel-generating}
For $z \in \CC$ and $t \neq 0$,
\[
\exp\Bigl(\frac z2\Bigl(t - \frac1t\Bigr)\Bigr) = \sum_{n \in \ZZ}J_n(z)t^n.
\]
Consequently, for integers $n$, $J_n(z) = \frac1\pi\int_0^\pi\cos(n\theta - z\sin\theta)\,d\theta$ (\emph{Bessel's integral}), and $|J_n(x)| \leq 1$ for real $x$.
\end{theorem}

\begin{proof}
Multiply the absolutely convergent series $e^{zt/2} = \sum_j\frac{(z/2)^jt^j}{j!}$ and $e^{-z/(2t)} = \sum_k\frac{(-1)^k(z/2)^kt^{-k}}{k!}$; the coefficient of $t^n$ ($n \geq 0$) is $\sum_k\frac{(-1)^k(z/2)^{n+2k}}{(n + k)!\,k!} = J_n(z)$,
and for $n < 0$ the analogous computation gives $(-1)^nJ_{-n} = J_n$. With $t = e^{i\theta}$, $e^{iz\sin\theta} = \sum_nJ_n(z)e^{in\theta}$, uniformly in $\theta$, so $J_n(z) = \frac{1}{2\pi}\int_{-\pi}^{\pi}e^{i(z\sin\theta - n\theta)}d\theta$; the
imaginary part of the integrand is odd in $\theta$, which gives Bessel's integral and the bound.
\end{proof}

\begin{proposition}\label{special-functions:proposition-bessel-recurrences}
$J_{\nu-1}(z) + J_{\nu+1}(z) = \frac{2\nu}{z}J_\nu(z)$ and $J_{\nu-1}(z) - J_{\nu+1}(z) = 2J_\nu^{\prime}(z)$. Moreover
$J_{1/2}(x) = \sqrt{\frac{2}{\pi x}}\sin x$ and $J_{-1/2}(x) = \sqrt{\frac{2}{\pi x}}\cos x$ for $x > 0$.
\end{proposition}

\begin{proof}
From the series, $(z^\nu J_\nu)^{\prime} = z^\nu J_{\nu-1}$ and $(z^{-\nu}J_\nu)^{\prime} = -z^{-\nu}J_{\nu+1}$, by differentiating termwise and using $\Gamma(k + \nu + 1) = (k + \nu)\Gamma(k + \nu)$. Expanding the derivatives gives
$J_\nu^{\prime} + \frac\nu zJ_\nu = J_{\nu-1}$ and $J_\nu^{\prime} - \frac\nu zJ_\nu = -J_{\nu+1}$; adding and subtracting gives the recurrences. For $\nu = \frac12$, the duplication formula gives
$k!\,\Gamma(k + \frac32) = \frac{(2k + 1)!\sqrt{\pi}}{2^{2k+1}}$, so $J_{1/2}(x) = \sqrt{\frac{2}{\pi x}}\sum_k\frac{(-1)^kx^{2k+1}}{(2k + 1)!}$; similarly for $\nu = -\frac12$.
\end{proof}

\begin{remark}\label{special-functions:remark-bessel-asymptotics}
Stationary phase (Remark~\ref{special-functions:remark-stationary-phase}) applied to Bessel's integral gives, for real $x \to \infty$,
$J_\nu(x) = \sqrt{\frac{2}{\pi x}}\bigl(\cos(x - \frac{\nu\pi}{2} - \frac\pi4) + O(x^{-1})\bigr)$, exact for $\nu = \pm\frac12$ by Proposition~\ref{special-functions:proposition-bessel-recurrences} \cite{watson-bessel}. Hence $J_\nu$ has
infinitely many positive zeros $j_{\nu,1} < j_{\nu,2} < \cdots$, spaced asymptotically by $\pi$. They are the frequencies of a circular drum: the eigenfunctions of the Laplacian on the unit disc with Dirichlet boundary
conditions are $J_n(j_{n,k}r)e^{\pm in\theta}$, with eigenvalues $j_{n,k}^2$, which one finds by separating variables in polar coordinates.
\end{remark}

Separating variables in spherical coordinates leads instead to spherical harmonics, which have a purely algebraic description.

\begin{definition}\label{special-functions:definition-harmonic-polynomials}
Let $\mathcal{P}_k$ be the space of homogeneous polynomials of degree $k$ in $n$ real variables, $|x|^2 = x_1^2 + \cdots + x_n^2$ and $\Delta = \sum_i\partial_i^2$. The space of \emph{harmonic polynomials of degree $k$}
is $\mathcal{H}_k = \{P \in \mathcal{P}_k : \Delta P = 0\}$; restrictions of its elements to $S^{n-1}$ are the \emph{spherical harmonics} of degree $k$.
\end{definition}

\begin{theorem}\label{special-functions:theorem-harmonic-decomposition}
$\mathcal{P}_k = \mathcal{H}_k \oplus |x|^2\mathcal{P}_{k-2}$, so $\mathcal{P}_k = \mathcal{H}_k \oplus |x|^2\mathcal{H}_{k-2} \oplus |x|^4\mathcal{H}_{k-4} \oplus \cdots$, and
\[
\dim\mathcal{H}_k = \binom{n + k - 1}{k} - \binom{n + k - 3}{k - 2}.
\]
For $n = 3$ this is $2k + 1$. Every polynomial restricted to $S^{n-1}$ is a sum of spherical harmonics.
\end{theorem}

\begin{proof}
\emph{An inner product.} For $P, Q \in \mathcal{P}_k$ put $\langle P, Q\rangle = P(\partial)Q$, the constant obtained by substituting $\partial_i$ for $x_i$ in $P$ and applying the operator to $Q$. On monomials,
$\langle x^\alpha, x^\beta\rangle = \alpha!\,\delta_{\alpha\beta}$, so this is a positive definite symmetric bilinear form. \emph{Adjointness.} For $R \in \mathcal{P}_{k-2}$, $\langle|x|^2R, Q\rangle = R(\partial)\Delta Q = \langle R, \Delta Q\rangle$,
since $|x|^2(\partial) = \Delta$. So $\mathcal{H}_k = \ker\Delta$ is the orthogonal complement of $|x|^2\mathcal{P}_{k-2}$, which gives the decomposition, and $\Delta \colon \mathcal{P}_k \to \mathcal{P}_{k-2}$ is surjective, being adjoint to
the injective multiplication by $|x|^2$. The dimension formula follows from $\dim\mathcal{P}_k = \binom{n + k - 1}{k}$ (with $\mathcal{P}_{-1} = \mathcal{P}_{-2} = 0$), and iterating gives the second decomposition. On the sphere $|x|^2 = 1$, so a polynomial of degree $k$
restricts to a sum of spherical harmonics of degrees $k, k - 2, \ldots$.
\end{proof}

\begin{remark}\label{special-functions:remark-spherical-harmonics}
The spherical harmonics of degree $k$ are eigenfunctions of the Laplacian of the round sphere with eigenvalue $-k(k + n - 2)$ (this is proved in Chapter~\ref{riemannian-geometry}), the
spaces $\mathcal{H}_k|_{S^{n-1}}$ are orthogonal in $L^2(S^{n-1})$, and by Theorem~\ref{special-functions:theorem-harmonic-decomposition} and the Stone--Weierstrass theorem their sum is dense: $L^2(S^{n-1}) = \widehat{\bigoplus}_k\mathcal{H}_k$.
For $n = 3$, the harmonics invariant under rotations about the $x_3$-axis are multiples of $P_k(\cos\theta)$, $\theta$ the polar angle, and the \emph{addition theorem} expresses $P_k(\langle\xi, \eta\rangle)$ as
$\frac{4\pi}{2k + 1}\sum_mY_{k,m}(\xi)\overline{Y_{k,m}(\eta)}$ for an orthonormal basis $(Y_{k,m})$ of $\mathcal{H}_k$. In quantum mechanics $\mathcal{H}_k$ is the space of states of angular momentum $k$, the irreducible representation of
$\mathrm{SO}(3)$ of dimension $2k + 1$.
\end{remark}

\section{Elliptic integrals and elliptic functions}\label{special-functions:section-elliptic}

\begin{definition}\label{special-functions:definition-elliptic-integrals}
For $0 \leq k < 1$, the \emph{complete elliptic integrals} of the first and second kind are
\[
K(k) = \int_0^{\pi/2}\frac{d\theta}{\sqrt{1 - k^2\sin^2\theta}}, \qquad E(k) = \int_0^{\pi/2}\sqrt{1 - k^2\sin^2\theta}\,d\theta,
\]
and $k^{\prime} = \sqrt{1 - k^2}$ is the complementary modulus.
\end{definition}

\begin{example}\label{special-functions:example-elliptic-integrals}
\begin{enumerate}
\item The ellipse $x = a\cos\theta$, $y = b\sin\theta$ with $a \geq b$ and eccentricity $e = \sqrt{1 - b^2/a^2}$ has perimeter $\int_0^{2\pi}\sqrt{a^2\sin^2\theta + b^2\cos^2\theta}\,d\theta = 4aE(e)$.
\item A pendulum of length $L$ released at angle $\theta_0$ satisfies $\ddot\theta = -\frac gL\sin\theta$. Energy conservation gives $\dot\theta^2 = \frac{2g}{L}(\cos\theta - \cos\theta_0)$, and the substitution
$\sin\frac\theta2 = \sin\frac{\theta_0}{2}\sin\phi$ turns the period $4\int_0^{\theta_0}d\theta/\dot\theta$ into $T = 4\sqrt{L/g}\,K(\sin\frac{\theta_0}{2})$. For small $\theta_0$, $K \approx \frac\pi2$ and $T \approx 2\pi\sqrt{L/g}$.
\end{enumerate}
\end{example}

\begin{definition}\label{special-functions:definition-agm}
For $a \geq b > 0$ let $a_0 = a$, $b_0 = b$, $a_{n+1} = \frac{a_n + b_n}{2}$, $b_{n+1} = \sqrt{a_nb_n}$. The \emph{arithmetic-geometric mean} $M(a, b)$ is the common limit of $(a_n)$ and $(b_n)$.
\end{definition}

The limit exists: $b_n \leq b_{n+1} \leq a_{n+1} \leq a_n$ by the inequality between the arithmetic and geometric means, and $a_{n+1} - b_{n+1} = \frac{(a_n - b_n)^2}{2(\sqrt{a_n} + \sqrt{b_n})^2}$, so
$a_{n+1} - b_{n+1} \leq \frac{(a_n - b_n)^2}{8b}$: the convergence is quadratic, and the number of correct digits roughly doubles at each step.

\begin{theorem}[Gauss]\label{special-functions:theorem-agm}
For $a, b > 0$,
\[
I(a, b) = \int_0^{\pi/2}\frac{d\theta}{\sqrt{a^2\cos^2\theta + b^2\sin^2\theta}} = \frac{\pi}{2M(a, b)}.
\]
In particular $K(k) = \frac{\pi}{2M(1, k^{\prime})}$.
\end{theorem}

\begin{proof}
\emph{Step 1.} The substitution $t = b\tan\theta$ gives $I(a, b) = \int_0^\infty\frac{dt}{\sqrt{(a^2 + t^2)(b^2 + t^2)}} = \frac12\int_{-\infty}^\infty\frac{dt}{\sqrt{(a^2 + t^2)(b^2 + t^2)}}$.
\emph{Step 2: invariance.} Let $a_1 = \frac{a + b}{2}$, $b_1 = \sqrt{ab}$, and substitute $t = \frac12(u - \frac{ab}{u})$, an increasing bijection from $(0, \infty)$ onto $\RR$ with $dt = \frac{u^2 + ab}{2u^2}du$. A direct computation gives
$a_1^2 + t^2 = \frac{(u^2 + a^2)(u^2 + b^2)}{4u^2}$ and $b_1^2 + t^2 = \frac{(u^2 + ab)^2}{4u^2}$, so $\frac{dt}{\sqrt{(a_1^2 + t^2)(b_1^2 + t^2)}} = \frac{2\,du}{\sqrt{(u^2 + a^2)(u^2 + b^2)}}$. By Step 1,
$I(a_1, b_1) = \frac12\cdot 2\int_0^\infty\frac{du}{\sqrt{(u^2 + a^2)(u^2 + b^2)}} = I(a, b)$.
\emph{Step 3.} Hence $I(a, b) = I(a_n, b_n)$ for all $n$, and $I$ is continuous, so $I(a, b) = I(M, M) = \frac{\pi}{2M}$ with $M = M(a, b)$. For $K$: $I(1, k^{\prime}) = \int_0^{\pi/2}\frac{d\theta}{\sqrt{1 - k^2\sin^2\theta}}$.
\end{proof}

\begin{example}\label{special-functions:example-lemniscate}
Gauss computed on 30 May 1799 that $M(1, \sqrt2) = 1.19814\ldots$ agrees to eleven decimals with $\pi/\varpi$, where $\varpi = 2\int_0^1\frac{dt}{\sqrt{1 - t^4}} = 2.62205\ldots$ is the \emph{lemniscate constant}, half the
length of the lemniscate $r^2 = \cos 2\theta$. Theorem~\ref{special-functions:theorem-agm} proves it: substituting $t = \cos\phi$, $\int_0^1\frac{dt}{\sqrt{1 - t^4}} = \int_0^{\pi/2}\frac{d\phi}{\sqrt{1 + \cos^2\phi}} = I(\sqrt2, 1) = \frac{\pi}{2M(\sqrt2, 1)}$.
The quadratic convergence of the AGM is the basis of the fastest known algorithms for computing $\pi$ (Brent--Salamin) \cite{borwein-pi-agm}.
\end{example}

The inverse functions of elliptic integrals are doubly periodic, as $\arcsin$ inverts $\int dt/\sqrt{1 - t^2}$ and $\sin$ has a period. Weierstrass constructed them directly.

\begin{definition}\label{special-functions:definition-lattice}
A \emph{lattice} in $\CC$ is $\Lambda = \ZZ\omega_1 + \ZZ\omega_2$ with $\omega_1$, $\omega_2$ linearly independent over $\RR$. An \emph{elliptic function} for $\Lambda$ is a meromorphic function $f$ on $\CC$ with $f(z + \omega) = f(z)$ for
all $\omega \in \Lambda$. The \emph{Eisenstein series} are $G_{2k}(\Lambda) = \sum_{\omega \in \Lambda \setminus 0}\omega^{-2k}$ for $k \geq 2$, and the \emph{Weierstrass function} is
\[
\wp(z) = \frac{1}{z^2} + \sum_{\omega \in \Lambda \setminus 0}\Bigl(\frac{1}{(z - \omega)^2} - \frac{1}{\omega^2}\Bigr).
\]
\end{definition}

\begin{lemma}\label{special-functions:lemma-lattice-sums}
$\sum_{\omega \in \Lambda \setminus 0}|\omega|^{-s}$ converges for $s > 2$. Hence the Eisenstein series converge absolutely, and the series for $\wp$ converges absolutely and uniformly on compact sets avoiding $\Lambda$.
\end{lemma}

\begin{proof}
Let $c > 0$ be the minimum of $|x\omega_1 + y\omega_2|$ on the square $\max(|x|, |y|) = 1$. The $8m$ lattice points with $\max(|x|, |y|) = m$ have $|\omega| \geq cm$, so the sum is at most $\sum_m8m(cm)^{-s} < \infty$. For $\wp$, when
$|z| \leq R$ and $|\omega| \geq 2R$, $\bigl|\frac{1}{(z - \omega)^2} - \frac{1}{\omega^2}\bigr| = \frac{|z||2\omega - z|}{|\omega|^2|z - \omega|^2} \leq \frac{10R}{|\omega|^3}$, and the case $s = 3$ applies.
\end{proof}

\begin{theorem}\label{special-functions:theorem-weierstrass-p}
$\wp$ is an even elliptic function with double poles exactly at $\Lambda$, and
\[
\wp(z) = \frac{1}{z^2} + \sum_{k \geq 1}(2k + 1)G_{2k+2}z^{2k} \quad \text{near } 0, \qquad (\wp^{\prime})^2 = 4\wp^3 - g_2\wp - g_3,
\]
with $g_2 = 60G_4$ and $g_3 = 140G_6$. A nonconstant elliptic function has poles; one without poles is constant.
\end{theorem}

\begin{proof}
\emph{Liouville.} An elliptic function without poles is bounded on the compact fundamental parallelogram, hence on $\CC$, hence constant (Corollary~\ref{complex-analysis:corollary-liouville}).
\emph{Periodicity.} By Lemma~\ref{special-functions:lemma-lattice-sums}, $\wp$ is meromorphic with double poles at $\Lambda$, and even since $\Lambda = -\Lambda$. Its derivative $\wp^{\prime}(z) = -2\sum_{\omega \in \Lambda}(z - \omega)^{-3}$ is
visibly periodic, so $\wp(z + \omega_i) - \wp(z) = C_i$ is constant; at $z = -\omega_i/2$ evenness gives $C_i = \wp(\omega_i/2) - \wp(-\omega_i/2) = 0$.
\emph{Laurent expansion.} For $|z| < |\omega|$, $\frac{1}{(z - \omega)^2} - \frac{1}{\omega^2} = \sum_{m \geq 1}(m + 1)\frac{z^m}{\omega^{m+2}}$; summing over $\omega$, the odd powers cancel by symmetry and the even ones give the
expansion.
\emph{Differential equation.} From $\wp = z^{-2} + 3G_4z^2 + 5G_6z^4 + O(z^6)$ one computes $(\wp^{\prime})^2 = 4z^{-6} - 24G_4z^{-2} - 80G_6 + O(z^2)$ and $4\wp^3 = 4z^{-6} + 36G_4z^{-2} + 60G_6 + O(z^2)$, so
$(\wp^{\prime})^2 - 4\wp^3 + 60G_4\wp + 140G_6$ is an elliptic function holomorphic at $0$, hence everywhere, vanishing at $0$; by Liouville it is identically $0$.
\end{proof}

\begin{remark}\label{special-functions:remark-elliptic-curves}
So $z \mapsto (\wp(z), \wp^{\prime}(z))$ maps $\CC/\Lambda$ minus $0$ onto the affine curve $y^2 = 4x^3 - g_2x - g_3$, and $0$ to the point at infinity: a complex torus is a plane cubic curve, an \emph{elliptic curve}. The
\emph{addition theorem} $\wp(z + w) = -\wp(z) - \wp(w) + \frac14\bigl(\frac{\wp^{\prime}(z) - \wp^{\prime}(w)}{\wp(z) - \wp(w)}\bigr)^2$ is the chord-and-tangent group law on the cubic. The coefficients $G_{2k}$, as functions of
$\tau = \omega_2/\omega_1$, are modular forms, and every elliptic function for $\Lambda$ is a rational function of $\wp$ and $\wp^{\prime}$. These themes are developed in the chapters on Riemann surfaces, elliptic curves and
modular forms.
\end{remark}

\section{Theta functions}\label{special-functions:section-theta}

\begin{definition}\label{special-functions:definition-theta}
For $z \in \CC$ and $\tau$ in the upper half-plane $\mathbb{H}$, the \emph{Jacobi theta function} is
\[
\vartheta(z \mid \tau) = \sum_{n \in \ZZ}e^{\pi in^2\tau + 2\pi inz}.
\]
\end{definition}

\begin{notation}\label{special-functions:notation-theta}
We write $q = e^{\pi i\tau}$, so $|q| < 1$, and $\vartheta(z \mid \tau) = \sum_nq^{n^2}e^{2\pi inz}$; $\vartheta(\tau) = \vartheta(0 \mid \tau)$. Here $\mathbb{H} = \{\operatorname{Im}\tau > 0\}$ is the upper half-plane, not the
quaternions.
\end{notation}

\begin{proposition}\label{special-functions:proposition-theta-basic}
The series converges absolutely and uniformly on compact subsets of $\CC \times \mathbb{H}$, so $\vartheta$ is holomorphic in both variables. It satisfies
\begin{enumerate}
\item $\vartheta(z + 1 \mid \tau) = \vartheta(z \mid \tau)$ and $\vartheta(z + \tau \mid \tau) = e^{-\pi i\tau - 2\pi iz}\vartheta(z \mid \tau)$;
\item the heat equation $\frac{\partial\vartheta}{\partial\tau} = \frac{1}{4\pi i}\frac{\partial^2\vartheta}{\partial z^2}$;
\item the product formula $\vartheta(z \mid \tau) = \prod_{n \geq 1}(1 - q^{2n})(1 + q^{2n-1}e^{2\pi iz})(1 + q^{2n-1}e^{-2\pi iz})$.
\end{enumerate}
Hence the zeros of $\vartheta(\cdot \mid \tau)$ are the points $\frac12 + \frac\tau2 + \ZZ + \ZZ\tau$, all simple.
\end{proposition}

\begin{proof}
For $\operatorname{Im}\tau \geq \varepsilon$ and $|\operatorname{Im}z| \leq R$, $|e^{\pi in^2\tau + 2\pi inz}| \leq e^{-\pi\varepsilon n^2 + 2\pi R|n|}$, which is summable; holomorphy follows from Weierstrass' theorem
(Corollary~\ref{complex-analysis:corollary-liouville}(4)), applied in each variable. (1) The first identity is clear; for the second, $\pi in^2\tau + 2\pi in(z + \tau) = \pi i(n + 1)^2\tau + 2\pi i(n + 1)z - \pi i\tau - 2\pi iz$, and reindexing gives the claim.
(2) Each term satisfies the equation, and termwise differentiation is allowed by locally uniform convergence. (3) This is Jacobi's triple product identity (Theorem~\ref{generating-functions:theorem-triple-product}) with
$z$ there replaced by $e^{2\pi iz}$: it is an identity of formal power series in $q$ whose coefficients are Laurent polynomials in $e^{2\pi iz}$, and for $|q| < 1$ both sides converge absolutely as double series, so the
coefficientwise identity becomes an identity of functions. The factors $1 + q^{2n-1}e^{\pm 2\pi iz}$ vanish exactly at the stated points (and $1 - q^{2n} \neq 0$), simply, and by Lemma~\ref{special-functions:lemma-infinite-products} these are all the zeros.
\end{proof}

To prove the transformation formula under $\tau \mapsto -1/\tau$ we need a version of the Poisson summation formula, which we prove with residues.

\begin{lemma}[Poisson summation in a strip]\label{special-functions:lemma-poisson-strip}
Let $f$ be holomorphic on a strip $|\operatorname{Im}z| < a$ with $|f(x + iy)| \leq \frac{A}{1 + x^2}$ there, and $\hat{f}(\xi) = \int_\RR f(x)e^{-2\pi ix\xi}dx$. Then $\sum_{n \in \ZZ}f(n) = \sum_{n \in \ZZ}\hat{f}(n)$, both series converging
absolutely.
\end{lemma}

\begin{proof}
Fix $0 < b < a$. The function $f(z)/(e^{2\pi iz} - 1)$ has simple poles at the integers with residues $f(n)/(2\pi i)$. Integrate it around the rectangle with vertices $\pm(N + \frac12) \pm ib$; on the vertical sides
$1/|e^{2\pi iz} - 1|$ is bounded independently of $N$ and $|f| = O(N^{-2})$, so by the residue theorem (Theorem~\ref{complex-analysis:theorem-residues}), letting $N \to \infty$,
\[
\sum_nf(n) = \int_{L_{-b}}\frac{f(z)}{e^{2\pi iz} - 1}dz - \int_{L_b}\frac{f(z)}{e^{2\pi iz} - 1}dz,
\]
where $L_{\pm b}$ is the line $\operatorname{Im}z = \pm b$ traversed from left to right. On $L_{-b}$, $|e^{2\pi iz}| = e^{2\pi b} > 1$ and $\frac{1}{e^{2\pi iz} - 1} = \sum_{n \geq 1}e^{-2\pi inz}$; on $L_b$, $|e^{2\pi iz}| = e^{-2\pi b} < 1$ and
$\frac{1}{e^{2\pi iz} - 1} = -\sum_{n \geq 0}e^{2\pi inz}$, both uniformly on the line. Integrating termwise, $\sum_nf(n) = \sum_{n \geq 1}\int_{L_{-b}}f(z)e^{-2\pi inz}dz + \sum_{n \geq 0}\int_{L_b}f(z)e^{2\pi inz}dz$. By Cauchy's theorem and the
decay of $f$, each line can be moved to the real axis, which gives $\sum_{n \geq 1}\hat{f}(n) + \sum_{n \geq 0}\hat{f}(-n)$. On $L_{\mp b}$ the terms are $O(e^{-2\pi nb})$, which gives the absolute convergence.
\end{proof}

\begin{lemma}\label{special-functions:lemma-gaussian-fourier}
For $t > 0$ and $\xi \in \RR$, $\int_\RR e^{-\pi tx^2}e^{-2\pi ix\xi}dx = t^{-1/2}e^{-\pi\xi^2/t}$.
\end{lemma}

\begin{proof}
By scaling it suffices to take $t = 1$. Completing the square, the integral is $e^{-\pi\xi^2}\int_\RR e^{-\pi(x + i\xi)^2}dx$. By Cauchy's theorem applied to $e^{-\pi z^2}$ on the rectangle with vertices $\pm R$, $\pm R + i\xi$, whose
vertical sides contribute $O(e^{-\pi(R^2 - \xi^2)})$, this equals $\int_\RR e^{-\pi x^2}dx = 1$ (Example~\ref{special-functions:example-gamma-half}).
\end{proof}

\begin{theorem}[Transformation formula]\label{special-functions:theorem-theta-transformation}
For $\tau \in \mathbb{H}$, $\vartheta(-1/\tau) = \sqrt{\tau/i}\,\vartheta(\tau)$, with the branch of the square root that is positive for $\tau \in i\RR_{>0}$. In particular, $\sum_ne^{-\pi n^2/t} = \sqrt{t}\sum_ne^{-\pi n^2t}$ for $t > 0$.
\end{theorem}

\begin{proof}
For $\tau = it$, apply Lemma~\ref{special-functions:lemma-poisson-strip} to $f(x) = e^{-\pi tx^2}$, which satisfies its hypotheses on every strip, since $|f(x + iy)| = e^{-\pi t(x^2 - y^2)}$; by Lemma~\ref{special-functions:lemma-gaussian-fourier},
$\sum_ne^{-\pi tn^2} = t^{-1/2}\sum_ne^{-\pi n^2/t}$, which is the claim for $\tau = it$. Both sides of the formula are holomorphic on $\mathbb{H}$ and agree on the imaginary axis, so they agree everywhere by the identity theorem
(Theorem~\ref{complex-analysis:theorem-identity}).
\end{proof}

\begin{remark}\label{special-functions:remark-theta-applications}
Together with $\vartheta(\tau + 2) = \vartheta(\tau)$, the transformation formula says that $\vartheta(\tau)^k$ is a modular form of weight $k/2$ for a congruence subgroup. Since $\vartheta(\tau)^k = \sum_nr_k(n)q^n$, where $r_k(n)$ is
the number of representations of $n$ as a sum of $k$ squares, the theory of modular forms gives formulas such as Jacobi's $r_4(n) = 8\sum_{d \mid n,\ 4 \nmid d}d$, which in particular reproves Lagrange's four-square theorem, and
$r_2(n) = 4\sum_{d \mid n}\chi_{-4}(d)$ \cite{hardy-wright}. The case $t$ real of Theorem~\ref{special-functions:theorem-theta-transformation} gives Riemann's functional equation of $\zeta$ in Chapter~\ref{analytic-number-theory}.
Physically, $\sum_ne^{-(x + n)^2/4t}/\sqrt{4\pi t}$ is the heat kernel on the circle $\RR/\ZZ$, and the transformation formula says that it equals the Fourier series $\sum_ne^{-4\pi^2n^2t + 2\pi inx}$. Quotients of theta functions with
characteristics give Jacobi's elliptic functions $\operatorname{sn}$, $\operatorname{cn}$, $\operatorname{dn}$, which invert the incomplete elliptic integral $\int_0^x\frac{dt}{\sqrt{(1 - t^2)(1 - k^2t^2)}}$.
\end{remark}

\section{Exercises}\label{special-functions:section-exercises}

\begin{exercise}\label{special-functions:exercise-gamma-half-integers}
Show that $\Gamma(n + \frac12) = \frac{(2n)!}{4^nn!}\sqrt{\pi}$ for $n \geq 0$, and that $\int_{-1}^1(1 - x^2)^ndx = \frac{2^{2n+1}(n!)^2}{(2n + 1)!}$.
\end{exercise}

\begin{solution}
By induction from $\Gamma(\frac12) = \sqrt\pi$ (Example~\ref{special-functions:example-gamma-half}) and $\Gamma(n + \frac32) = (n + \frac12)\Gamma(n + \frac12)$: $\frac{(2n + 2)!}{4^{n+1}(n + 1)!} = \frac{(2n + 1)(2n + 2)}{4(n + 1)}\cdot\frac{(2n)!}{4^nn!}$ and
$\frac{(2n + 1)(2n + 2)}{4(n + 1)} = n + \frac12$. For the integral, substitute $u = x^2$ on $[0, 1]$: $2\int_0^1(1 - x^2)^ndx = \int_0^1u^{-1/2}(1 - u)^ndu = B(\frac12, n + 1) = \frac{\Gamma(\frac12)n!}{\Gamma(n + \frac32)}$, and the first
part gives $\frac{\sqrt\pi\,n!\,4^{n+1}(n + 1)!}{(2n + 2)!\sqrt\pi} = \frac{2^{2n+1}(n!)^2}{(2n + 1)!}$.
\end{solution}

\begin{exercise}\label{special-functions:exercise-duplication-bohr-mollerup}
Use Theorem~\ref{special-functions:theorem-bohr-mollerup} to prove the duplication formula for real $x > 0$.
\end{exercise}

\begin{solution}
Let $f(x) = \frac{2^{x-1}}{\sqrt\pi}\Gamma(\frac x2)\Gamma(\frac{x + 1}{2})$. Its logarithm is a sum of a linear function and two convex functions, so it is convex. $f(x + 1) = \frac{2^x}{\sqrt\pi}\Gamma(\frac{x+1}2)\Gamma(\frac x2 + 1) = 2\cdot\frac x2\cdot f(x) = xf(x)$,
and $f(1) = \frac{1}{\sqrt\pi}\Gamma(\frac12)\Gamma(1) = 1$. So $f = \Gamma$, which is $\Gamma(\frac x2)\Gamma(\frac x2 + \frac12) = 2^{1-x}\sqrt\pi\,\Gamma(x)$, the duplication formula with $s = x/2$.
\end{solution}

\begin{exercise}\label{special-functions:exercise-digamma}
Show that $\Gamma^{\prime}(1) = -\gamma$ and $\frac{\Gamma^{\prime}}{\Gamma}(s) = -\gamma - \frac1s + \sum_{k \geq 1}\bigl(\frac1k - \frac{1}{s + k}\bigr)$.
\end{exercise}

\begin{solution}
Take the logarithmic derivative of the Weierstrass product (Theorem~\ref{special-functions:theorem-weierstrass-gamma}) with Lemma~\ref{special-functions:lemma-infinite-products}:
$-\frac{\Gamma^{\prime}}{\Gamma}(s) = \frac1s + \gamma + \sum_k\bigl(\frac{1}{s + k} - \frac1k\bigr)$. At $s = 1$, $\sum_k(\frac{1}{k+1} - \frac1k) = -1$, so $-\Gamma^{\prime}(1) = 1 + \gamma - 1 = \gamma$, using $\Gamma(1) = 1$.
\end{solution}

\begin{exercise}\label{special-functions:exercise-watson-example}
Find the asymptotic expansion of $\int_0^\infty\frac{e^{-xt}}{1 + t^2}dt$ as $x \to \infty$.
\end{exercise}

\begin{solution}
$\frac{1}{1 + t^2} = \sum_k(-1)^kt^{2k}$ near $0$, so Watson's lemma (Theorem~\ref{special-functions:theorem-watson}, $\alpha = 1$, $a_{2k} = (-1)^k$, odd coefficients $0$) gives $\sum_k(-1)^k\frac{(2k)!}{x^{2k+1}} = \frac1x - \frac{2}{x^3} + \frac{24}{x^5} - \cdots$.
\end{solution}

\begin{exercise}\label{special-functions:exercise-gauss-legendre}
Show that the two-point Gauss rule for $w = 1$ on $(-1, 1)$ is $\int_{-1}^1q = q(-\frac{1}{\sqrt3}) + q(\frac{1}{\sqrt3})$, exact for cubic polynomials, and that it fails for $q = x^4$.
\end{exercise}

\begin{solution}
The monic orthogonal polynomial of degree $2$ is $x^2 - \frac13$ (orthogonal to $1$ and $x$), with zeros $\pm\frac{1}{\sqrt3}$. By symmetry the two weights are equal, and $\int 1 = 2$ makes them $1$. By
Theorem~\ref{special-functions:theorem-gauss-quadrature} the rule is exact on $\mathcal{P}_3$. For $x^4$: $\int_{-1}^1x^4 = \frac25$, while the rule gives $\frac29$.
\end{solution}

\begin{exercise}\label{special-functions:exercise-legendre-hypergeometric}
Show that $F(-n, n + 1; 1; \frac{1 - x}{2})$ is a polynomial of degree $n$ in $x$ with value $1$ at $x = 1$, solving Legendre's equation with $\lambda = n(n + 1)$, and conclude that it equals $P_n(x)$.
\end{exercise}

\begin{solution}
\emph{The hypergeometric polynomial.} The series terminates since $(-n)_k = 0$ for $k > n$, and the coefficient of $u^n$ is nonzero, so it is a polynomial of degree $n$ in $u = \frac{1 - x}{2}$, hence in $x$;
at $x = 1$, $u = 0$ and $F = 1$. By Proposition~\ref{special-functions:proposition-hypergeometric-equation} with $(a, b, c) = (-n, n + 1, 1)$ it solves $u(1 - u)w^{\prime\prime} + (1 - 2u)w^{\prime} + n(n + 1)w = 0$, and since
$u(1 - u) = \frac{1 - x^2}{4}$, $\frac{d}{du} = -2\frac{d}{dx}$ and $1 - 2u = x$, this becomes Legendre's equation $(1 - x^2)w^{\prime\prime} - 2xw^{\prime} + n(n + 1)w = 0$.
\emph{$P_n$ solves it.} Let $L(y) = ((1 - x^2)y^{\prime})^{\prime} + n(n + 1)y$. Then $L(P_n)$ has degree at most $n$, and its coefficient of $x^n$ is $(-n(n + 1) + n(n + 1))$ times that of $P_n$, so $\deg L(P_n) < n$. For
$q \in \mathcal{P}_{n-1}$, integrating by parts twice (the boundary terms contain $1 - x^2$), $\int_{-1}^1qL(P_n) = \int_{-1}^1L(q)P_n = 0$, by Proposition~\ref{special-functions:proposition-legendre-hermite}(1). So $L(P_n)$ is orthogonal to
itself, and $L(P_n) = 0$.
\emph{Uniqueness.} For $y = \sum_jc_jx^j$, $L(y) = 0$ means $(j + 2)(j + 1)c_{j+2} = (j(j + 1) - n(n + 1))c_j$ for all $j$. If $y$ is a polynomial solution, let $j$ be the largest index with $c_j \neq 0$; the relation gives
$j(j + 1) = n(n + 1)$, so $j = n$. Going down from $c_n$, the relation determines $c_{n-2}, c_{n-4}, \ldots$, since $j(j + 1) - n(n + 1) \neq 0$ for $j < n$; and the coefficients of the other parity vanish, since their
largest nonzero one would again have index $n$. So the polynomial solutions form a line, and the two polynomials, both equal to $1$ at $x = 1$, coincide.
\end{solution}

\begin{exercise}\label{special-functions:exercise-bessel-integral}
Show that $J_0^{\prime} = -J_1$ and $\int_0^xtJ_0(t)\,dt = xJ_1(x)$.
\end{exercise}

\begin{solution}
By Proposition~\ref{special-functions:proposition-bessel-recurrences} with $\nu = 0$: $(z^0J_0)^{\prime} = -z^0J_1$, that is $J_0^{\prime} = -J_1$; with $\nu = 1$: $(zJ_1)^{\prime} = zJ_0$, and integrating from $0$, where $zJ_1(z)$ vanishes,
gives the second identity.
\end{solution}

\begin{exercise}\label{special-functions:exercise-harmonic-degree-two}
Find a basis of $\mathcal{H}_2$ for $n = 3$ and check $\dim\mathcal{H}_2 = 5$.
\end{exercise}

\begin{solution}
$\mathcal{P}_2$ has dimension $6$, and $\Delta$ maps it onto the constants, with $\Delta(\sum a_{ij}x_ix_j) = 2\sum_ia_{ii}$. So $\mathcal{H}_2$ consists of the quadratic forms with traceless matrix: a basis is
$x_1x_2$, $x_1x_3$, $x_2x_3$, $x_1^2 - x_2^2$, $x_2^2 - x_3^2$, and $5 = \binom42 - \binom20$, in agreement with Theorem~\ref{special-functions:theorem-harmonic-decomposition}.
\end{solution}

\begin{exercise}\label{special-functions:exercise-agm-numerics}
Compute three steps of the AGM of $1$ and $\sqrt2$ and deduce $\varpi$ to five digits.
\end{exercise}

\begin{solution}
$(a, b) = (1.41421, 1)$, then $(1.207107, 1.189207)$, then $(1.198157, 1.198124)$, then $(1.1981402, 1.1981402)$. So $M(1, \sqrt2) = 1.19814$ and $\varpi = \pi/M = 2.62206$ (more precisely $2.622057\ldots$),
by Example~\ref{special-functions:example-lemniscate}.
\end{solution}

\begin{exercise}\label{special-functions:exercise-wp-half-periods}
Show that $\wp^{\prime}$ vanishes at the half-periods $\frac{\omega_1}{2}$, $\frac{\omega_2}{2}$, $\frac{\omega_1 + \omega_2}{2}$, so that $4x^3 - g_2x - g_3 = 4(x - e_1)(x - e_2)(x - e_3)$ with $e_i$ the values of $\wp$ there.
\end{exercise}

\begin{solution}
$\wp^{\prime}$ is odd and periodic, so for a half-period $h$, $\wp^{\prime}(h) = -\wp^{\prime}(-h) = -\wp^{\prime}(h - 2h) = -\wp^{\prime}(h)$, hence $\wp^{\prime}(h) = 0$ (the half-periods are not poles). By the differential equation each
$e_i = \wp(h_i)$ is a root of $4x^3 - g_2x - g_3$. They are distinct: $\wp - e_i$ has a double zero at $h_i$ (as $\wp^{\prime}(h_i) = 0$), and an elliptic function of order $2$ (two poles counted with multiplicity in a period
parallelogram) has exactly two zeros there, by the argument principle (Theorem~\ref{complex-analysis:theorem-argument-principle}) applied to the boundary of the parallelogram, where the contributions of opposite sides cancel.
So $\wp(h_j) \neq e_i$ for $j \neq i$.
\end{solution}
